\SourcePage{303}
\SGAsection{7}{光滑群的 Cartan 子群}
\begin{SGAresult}{命题}{7.1}
设 $S$ 是\Term{prescheme}{预概形}，$G$ 是一个在 $S$ 上\Term{finite-presentation}{有限呈示}、\Term{smooth}{光滑}且\Term{fiber}{纤维}\Term{connected}{连通}的 $S$-\Term{group-prescheme}{群预概形}。
\SourcePage{304}
\begin{enumerate}
\item[i)] 设 $\mathcal{CT}$ 是如下 $S$-\Term{functor}{函子} $(\mathrm{Sch}/S)^\circ\to\mathrm{Ens}$：对于每个 $S$-\Term{prescheme}{预概形} $S'$，$\mathcal{CT}(S')$ 是 $G_{S'}$ 中所有在 $S'$ 上\Term{smooth}{光滑}的\Term{group-subprescheme}{群子预概形} $H$ 所成的\Term{set}{集合}，并要求对 $S'$ 的每一点 $s'$，$H_{s'}$ 是 $G_{s'}$ 中某个\Term{subtorus}{子环面}的\Term{centralizer}{中心化子}。则 $\mathcal{CT}$ 可由一个在 $S$ 上\Term{finite-presentation}{有限呈示}、\Term{smooth}{光滑}且\Term{quasi-projective}{拟射影}的 $S$-\Term{prescheme}{预概形}表示。
\item[ii)] 若 $S$ 是 artin 的，且 $H\in\mathcal{CT}(S)$，则 $H$ 是 $G$ 中某个\Term{subtorus}{子环面}的\Term{centralizer}{中心化子}。若 $S$ 是 hensel \Term{local-ring}{局部环}的谱，且 $H\in\mathcal{CT}(S)$，则 $H$ 是 $G$ 中某个在 $S$ 上\Term{etale}{平展}的\Term{multiplicative-type}{乘法型}\Term{subgroup}{子群}的\Term{centralizer}{中心化子}。
\item[iii)] 若 $L$ 是在 $S$ 上\Term{smooth}{光滑}且\Term{finite-presentation}{有限呈示}的 $S$-\Term{group-prescheme}{群预概形}，$i:L\to G$ 是 $S$-群\Term{monomorphism}{单态射}，$H\in\mathcal{CT}(S)$，则 $\operatorname{Transp}_G(H,L)$ 和 $\operatorname{Transp}^{\mathrm{str}}_G(H,L)$（参见第 VIII 讲 6.5 e)）可由 $G$ 中在 $S$ 上\Term{smooth}{光滑}的闭\Term{subprescheme}{子预概形}表示。
\item[iv)] 若 $H\in\mathcal{CT}(S)$，则 $H$ 在 $G$ 中闭，$N=\operatorname{Norm}_G(H)$ 可由 $G$ 中在 $S$ 上闭且\Term{smooth}{光滑}的\Term{group-subprescheme}{群子预概形}表示；$N/H$ 可由一个在 $S$ 上\Term{separated}{分离}、\Term{etale}{平展}且\Term{finite-type}{有限型}的 $S$-\Term{group-prescheme}{群预概形}表示；$G/N$ 可由一个在 $S$ 上\Term{smooth}{光滑}且\Term{quasi-projective}{拟射影}的 $S$-\Term{prescheme}{预概形}表示。
\item[v)] 设 $G'$ 是在 $S$ 上\Term{finite-presentation}{有限呈示}的 $S$-\Term{group-prescheme}{群预概形}，$u:G\to G'$ 是\Term{faithfully-flat}{忠实平坦}的 $S$-\Term{group-morphism}{群态射}，因而 $G'$ 满足与 $G$ 相同的假设（第 VI$_B$ 讲 9）。若 $H\in\mathcal{CT}_G(S)$，则 $H$ 经 $u$ 的像可由 $G'$ 的一个\Term{group-subprescheme}{群子预概形} $H'$ 表示，且 $H'\in\mathcal{CT}_{G'}(S)$。此外，$H\to H'$ \Term{faithfully-flat}{忠实平坦}；若 $H$ 是 $G$ 中\Term{torus}{环面} $T$ 的\Term{centralizer}{中心化子}，则 $H'$ 是 $G'$ 中\Term{torus}{环面} $T'=u(T)$ 的\Term{centralizer}{中心化子}。
\item[vi)] 在 v) 的条件下，考虑 $S$-\Term{morphism}{态射}
\[\widetilde u:\mathcal{CT}_G\longrightarrow\mathcal{CT}_{G'},\qquad H\longmapsto H'=u(H).\]
则 $\widetilde u$ 是\Term{faithfully-flat}{忠实平坦}、\Term{quasi-compact}{拟紧}的\Term{morphism}{态射}；若再假定 $\operatorname{Ker}u$ 是\Term{central}{中心的}，则 $\widetilde u$ 是\Term{isomorphism}{同构}，其逆\Term{isomorphism}{同构}为 $H'\mapsto u^{-1}(H')$。
\end{enumerate}
\end{SGAresult}
\textbf{ii) 的证明。} 对第一条断言，可以假定 $S$ 是\Term{closed-point}{闭点}为 $s$ 的 artin \Term{local-scheme}{局部}\Term{prescheme}{预概形}。设 $H\in\mathcal{CT}(S)$，$T_s$ 为 $H_s$ 的极大\Term{center}{中心}\Term{torus}{环面}；它已在 $\kappa(s)$ 上定义（参见 3.4）。由于 $H_s\in\mathcal{CT}(s)$，有 $H_s=\operatorname{Centr}_{G_s}T_s$。群 $H$ \Term{smooth}{光滑}，所以 $T_s$ 可以（唯一地）\Term{lifting}{提升}为 $H$ 中的一个\Term{center}{中心}\Term{subtorus}{子环面} $T$（第 IX 讲 3.6 bis 和 5.6）。于是 $H'=\operatorname{Centr}_G T$ 包含 $H$，且与 $H$ 有相同的\Term{fiber}{纤维}；由于 $H$ 在 $S$ 上\Term{flat}{平坦}，有 $H=H'$（第 VI$_B$ 讲 2.5）。

现在设 $S$ 是 hensel \Term{local-ring}{局部环}的谱；通过通常的\Term{reduced}{约化}，可以假定它是 noether 的。以 $s$ 表示 $S$ 的\Term{closed-point}{闭点}，$T_s$ 表示 $H_s$ 的极大\Term{center}{中心}\Term{torus}{环面}，$q$ 为在 $S$ 上\Term{invertible}{可逆}的整数，$\ell$ 为 $q$ 的一个幂，$T$ 为闭\Term{fiber}{纤维}\Term{isomorphism}{同构}于 $T_s$ 的 $S$-\Term{torus}{环面}（第 X 讲 4.6）。另一方面，以 ${}_\ell H$ 表示 $H$ 中 $\ell$ 次幂映射的“核”，$U_\ell$ 表示 ${}_\ell H$ 中在 $S$ 上\Term{etale}{平展}的最大\Term{open-subset}{开集}。由 1.3 及 $H$ 在 $S$ 上\Term{flat}{平坦}可知，$U_\ell$ 包含 ${}_\ell T_s$。由于 $S$ 是 hensel 的，存在唯一的 $S$-\Term{morphism}{态射}
\[u_\ell:{}_\ell T\longrightarrow U_\ell\]
\SourcePage{305}
在闭\Term{fiber}{纤维}上诱导\Term{canonical}{典范}\Term{immersion}{浸入} ${}_\ell T_s\to(U_\ell)_s$。由唯一性，容易看出 $u_\ell$ 是一个\Term{center}{中心} $S$-\Term{group-morphism}{群态射}（第 IX 讲 5.6 a)）。仿照 6.6 和 6.7，可证明 $u_\ell$ 是\Term{immersion}{浸入}，而且若 $M_\ell=u_\ell({}_\ell T)$ 是像群，则当 $\ell$ 充分大时，$\operatorname{Centr}_G(M_\ell)=H$（此处用到 $S$ 是 noether 的假设）。

\textbf{i) 的证明。} 群 $G$ 在 $S$ 上\Term{smooth}{光滑}且\Term{fiber}{纤维}\Term{connected}{连通}，所以若 $H\in\mathcal{CT}(S)$，则 $H$ 的\Term{fiber}{纤维}\Term{connected}{连通}（第 XII 讲 6.6 b)），且 $H$ 等于其\Term{connected-normalizer}{连通正规化子}（6.5）；因此 7.1 i) 定义的\Term{functor}{函子} $\mathcal{CT}$ 与 6.4.0 定义的同名\Term{functor}{函子}一致。由\Term{theorem}{定理} 6.4，$\mathcal{CT}$ 可由在 $S$ 上\Term{finite-presentation}{有限呈示}且\Term{quasi-projective}{拟射影}的 $S$-\Term{prescheme}{预概形}表示。还须证明此\Term{prescheme}{预概形}在 $S$ 上\Term{smooth}{光滑}。先用 EGA IV 8 归结到 $S$ 是\Term{affine}{仿射} noether 的情形。由第 XI 讲 1.5，只须证明：若 $S$ 是 artin \Term{local-ring}{局部环}的谱，$S_0$ 是由\Term{nilpotent-ideal}{幂零理想}定义的\Term{subscheme}{子概形}，$H_0\in\mathcal{CT}(S_0)$，则 $H_0$ \Term{liftable}{可提升}为 $G$ 中在 $S$ 上\Term{smooth}{光滑}的\Term{group-subprescheme}{群子预概形} $H$。由 ii)，有 $H_0=\operatorname{Centr}_{G_0}T_0$，其中 $T_0$ 是 $G_0$ 的\Term{subtorus}{子环面}。由于 $G$ \Term{smooth}{光滑}，$T_0$ \Term{liftable}{可提升}为 $G$ 的\Term{subtorus}{子环面} $T$（第 IX 讲 3.6 bis）；取 $H=\operatorname{Centr}_G(T)$ 即可，它确在 $S$ 上\Term{smooth}{光滑}（第 XI 讲 2.4 及\Term{lemma}{引理} 2.5）。

\textbf{iii) 的证明。} 由于 $H$ \Term{smooth}{光滑}且\Term{fiber}{纤维}\Term{connected}{连通}，由第 XI 讲 6.11，\Term{transporter}{输运子} $\operatorname{Transp}_G(H,L)$ 可由 $G$ 中在 $S$ 上\Term{finite-presentation}{有限呈示}的闭\Term{subprescheme}{子预概形}表示。为证明此\Term{transporter}{输运子}\Term{smooth}{光滑}，如上归结到以下\Term{lifting}{提升}问题：$S$ 是 artin \Term{local-scheme}{局部}的，$S_0$ 是 $S$ 的闭\Term{subprescheme}{子预概形}，$g_0\in G(S_0)$ 满足 $\operatorname{int}(g_0)H_0\subset L_0$，须将 $g_0$ \Term{lifting}{提升}为满足 $\operatorname{int}(g)H\subset L$ 的 $g\in G(S)$。由于 $G$ 在 $S$ 上\Term{smooth}{光滑}，存在\Term{lifting}{提升} $g_0$ 的 $G$ 的\Term{section}{截面} $g_1$；置 $H'=\operatorname{int}(g_1)H$。于是 $H'\in\mathcal{CT}(S)$，且 $H'_0\subset L_0$。由 ii)，$H'$ 是 $G$ 中某个\Term{torus}{环面} $T'$ 的\Term{centralizer}{中心化子}。由于 $L$ \Term{smooth}{光滑}，$L_0$ 的\Term{torus}{环面} $T'_0$ \Term{liftable}{可提升}为 $L$ 的\Term{torus}{环面} $T''$（第 XI 讲 3.6 bis）。群 $\operatorname{Centr}_L T''$ 包含于 $\operatorname{Centr}_G T''$，与后者有相同的\Term{fiber}{纤维}（即 $H'_0$），且\Term{smooth}{光滑}，故等于 $\operatorname{Centr}_G T''=H''$。$G$ 的\Term{subtorus}{子环面} $T'$ 和 $T''$ 是 $T'_0$ 的两个\Term{lifting}{提升}，因而由某个\Term{reduced}{约化}为 $G_0$ 的\Term{unit-section}{单位截面}的 $h\in G(S)$ \Term{conjugacy}{共轭}（第 IX 讲 3.3 bis）；它们在 $G$ 中的\Term{centralizer}{中心化子} $H'$ 和 $H''$ 也如此。\Term{section}{截面} $g=hg_1$ \Term{lifting}{提升} $g_0$，并确实满足 $\operatorname{int}(g)H\subset L$。

若现在 $g\in\operatorname{Transp}_G(H,L)(S)$，则 $\operatorname{int}(g)H=L$ 的充要条件是，对每个 $s\in S$ 都有 $\dim H_s=\dim L_s$。因此，以 $U$ 表示 $S$ 中使 $H$ 与 $L$ 的\Term{fiber}{纤维}\Term{dimension}{维数}相同的既开又闭\Term{subprescheme}{子预概形}，则 $H$ 到 $L$ 的\Term{strict-transporter}{严格输运子} $\operatorname{Transp}^{\mathrm{str}}_G(H,L)$ 可由如下 $S$-\Term{prescheme}{预概形}表示：
\[U\times_S\operatorname{Transp}_G(H,L).\]

\textbf{iv) 的证明。} 为证明 $H\in\mathcal{CT}(S)$ 时 $H$ 在 $G$ 中闭，可以先假定 $S$ \Term{affine}{仿射} noether，再假定 $S$ 是\Term{complete}{完备}\Term{local-ring}{局部环}的谱（EGA IV 8）。此时由 ii)，$H$ 是 $G$ 中某个\Term{multiplicative-type}{乘法型}\Term{subgroup}{子群}的\Term{centralizer}{中心化子}，而 $G$ 在 $S$ 上\Term{separated}{分离}（第 VI$_B$ 讲 5.2），所以 $H$ 是闭的。
\SourcePage{306}
由 iii)，$N=\operatorname{Norm}_G(H)=\operatorname{Transp}^{\mathrm{str}}_G(H,H)$ 可由 $G$ 中在 $S$ 上\Term{smooth}{光滑}的闭\Term{group-subprescheme}{群子预概形}表示。考虑 $S$-\Term{morphism}{态射}
\[G\longrightarrow\mathcal{CT},\qquad g\longmapsto\operatorname{int}(g)H.\]
由 iii)，此\Term{morphism}{态射}\Term{smooth}{光滑}，故其像是 $\mathcal{CT}$ 中的\Term{open-subset}{开集} $U$。仿照第 XI 讲 5.3，可证明 $G/N$ 由 $U$ 表示，故尤其是\Term{quasi-projective}{拟射影}的。

现在研究商 $N/H$。由 EGA IV 8，为证明 $N/H$ \Term{representable}{可表}，可以假定 $S$ \Term{affine}{仿射} noether，继而假定 $S$ 是\Term{local-ring}{局部环} $A$ 的谱，从而\Term{dimension}{维数}\footnote{编者注：（\Term{krull-dimension}{Krull 维数}）。}\Term{finite}{有限}。对 $S$ 的\Term{dimension}{维数}作升归纳。若 $\dim S=0$，结论由第 VI$_A$ 讲 \S4 得到。注意，若 $N/H$ \Term{representable}{可表}，则它在 $S$ 上\Term{separated}{分离}（因为由 iv)，$H$ 在 $N$ 中闭），且在 $S$ 上\Term{finite-type}{有限型}、\Term{etale}{平展}（因为 $N$ 在 $S$ 上\Term{smooth}{光滑}且\Term{finite-type}{有限型}，$H$ 在 $N$ 中开），所以 $N/H$ 必在 $S$ 上\Term{quasi-affine}{拟仿射}（SGA 1 VIII 6.2）。由\Term{quasi-affine}{拟仿射}\Term{scheme}{概形}的\Term{effective}{有效}\Term{descent}{下降}（同上，7.9），可以用 $A$ 的\Term{completion}{完备化}代替 $A$，从而假定 $S$ 是\Term{complete}{完备} noether \Term{local-ring}{局部环}的谱。

以 $s$ 表示其\Term{closed-point}{闭点}，$U=S\setminus\{s\}$。由归纳假设，$(N|_U)/(H|_U)$ 可由 $U$-群 $K$ 表示。设 $T_s$ 是 $H_s$ 的极大\Term{center}{中心}\Term{torus}{环面}，$q$ 是在 $S$ 上\Term{invertible}{可逆}的整数，$\ell$ 是 $q$ 的一个幂，$M_\ell$ 是 $H$ 中唯一\Term{lifting}{提升} ${}_\ell T_s$ 的\Term{center}{中心}\Term{multiplicative-type}{乘法型}\Term{subgroup}{子群}（参见 ii)）。选取充分大的 $\ell$，使 $\operatorname{Centr}_G(M_\ell)=H$，并令 $N'=\operatorname{Norm}_G(M_\ell)$。由 $\operatorname{Centr}_G(M_\ell)=H$，有 $N'\subset\operatorname{Norm}_G(H)$。此外，直接可见 $T_s$（从而 $(M_\ell)_s$）是 $H_s$ 的\Term{characteristic-subgroup}{特征子群}，即在 $\operatorname{Aut}_{\mathrm{gr}}(H_s)$ 下\Term{stable}{稳定}，故 $N_s$ \Term{normalization}{正规化} $(M_\ell)_s$，从而 $N_s=N'_s$。第 XI 讲 5.9 的证明说明，商
\[N'/H=\operatorname{Norm}_G(M_\ell)/\operatorname{Centr}_G(M_\ell)\]
可由 $S$-群 $K'$ 表示。由于 $N'$ 在 $S$ 上\Term{smooth}{光滑}且 $N'_s=N_s$，$N'$ 是 $N$ 的一个包含 $H$ 的\Term{open-subgroup}{开子群}（第 VI$_B$ 讲 2.5）；所以 $N'|_U$ 在 $K$ 中的像是一个\Term{isomorphism}{同构}于 $K'|_U$ 的\Term{open-subgroup}{开子群}。令 $L$ 为利用此\Term{isomorphism}{同构}将 $K$ 与 $K'$ \Term{gluing}{黏合}而得的 $S$-\Term{prescheme}{预概形}，$p:N\to L$ 为\Term{gluing}{黏合}\Term{canonical-projection}{典范投影} $N|_U\to K$ 和 $N'\to K'$ 而得的 $S$-\Term{morphism}{态射}。显然 $(L,p)$ 表示商 $N/H$。

\textbf{v) 的证明。} 先假定 $S$ 是域 $k$ 的谱。$H$ 的像 $H'$ 是 $G'$ 的\Term{smooth}{光滑}\Term{subgroup}{子群}。须证明 $H'\in\mathcal{CT}_{G'}(S)$，这将由下面更精确的\Term{lemma}{引理}得到。
\begin{SGAresult}{\Term{lemma}{引理}}{7.2}
设 $u:G\to G'$ 是\Term{smooth}{光滑}\Term{connected}{连通}\Term{algebra}{代数} $k$-群的\Term{epimorphism}{满态射}，$T$ 是 $G$ 的\Term{torus}{环面}，$T'$ 是它在 $G'$ 中的像，则
\[u(\operatorname{Centr}_G T)=\operatorname{Centr}_{G'}T'.\]
\end{SGAresult}
置 $H=\operatorname{Centr}_G T$，$H'=\operatorname{Centr}_{G'}T'$，$H''=u(H)$。有 $H''\subset H'$。为证明 $H'=H''$，可以假定\Term{base-field}{基域}\Term{algebraically-closed}{代数闭}，且只须证明 $H'$ 的每个\Term{cartan-subgroup}{Cartan 子群}都包含于 $H''$。事实上，此时 $H''$ 将包含 $H'$ 的\Term{regular-element}{正则点}\Term{open-subset}{开集}，故 $H''$ 是 $H'$ 的\Term{open-subgroup}{开子群}，
\SourcePage{307}
因后者的\Term{fiber}{纤维}\Term{connected}{连通}，便有 $H''=H'$。于是设 $C'$ 是 $H'$ 的\Term{cartan-subgroup}{Cartan 子群}；$C'$ 也是 $G'$ 的\Term{cartan-subgroup}{Cartan 子群}，因为 $H'$ 是\Term{torus}{环面} $T'$ 的\Term{centralizer}{中心化子}，从而与 $G'$ 有相同的\Term{reductive-rank}{约化秩}和\Term{nilpotent-rank}{幂零秩}。置 $K=(u^{-1}(C'))^0_{\mathrm{r\text{é}d}}$。由于 $T'$ 位于 $H'$ 的\Term{center}{中心}，$C'$ 包含 $T'$，故 $K$ 包含 $T$。取 $K$ 中包含 $T$ 的一个\Term{cartan-subgroup}{Cartan 子群} $C$。\Term{torus}{环面} $T$ 包含于 $C$ 的唯一\Term{maximal-torus}{极大环面}中，而后者在 $C$ 中是\Term{central}{中心的}（第 XII 讲 6.6 c)），故 $C$ 包含于 $H=\operatorname{Centr}_G T$。利用 $G$ 的任意两个\Term{cartan-subgroup}{Cartan 子群}\Term{conjugacy}{共轭}，以及 $G$ 的\Term{cartan-subgroup}{Cartan 子群}的像是 $G'$ 的\Term{cartan-subgroup}{Cartan 子群}（第 XII 讲 6.6），可知 $C$ 也是 $G$ 的\Term{cartan-subgroup}{Cartan 子群}；故其像是 $G'$ 的\Term{cartan-subgroup}{Cartan 子群}。此像包含于 $C'$，所以 $u(C)=C'$，从而 $C'$ 确实包含于 $H''$。

这就证明了 $S$ 是域 $k$ 的谱时的 v)。现在讨论一般情形。由于 $G,H,G'$ 在 $S$ 上\Term{finite-presentation}{有限呈示}，为证明 $u(H)$ \Term{representable}{可表}且属于 $\mathcal{CT}_{G'}(S)$，可以用通常的方法先归结到 $S$ \Term{affine}{仿射} noether，继而归结到 $S$ 是\Term{local-ring}{局部环}的谱。再由 $G'$ 的\Term{subprescheme}{子预概形}的 fpqc \Term{descent}{下降}，甚至可以假定 $S$ 是\Term{complete}{完备} noether \Term{local-ring}{局部环} $A$ 的谱。

沿用 ii) 的记号：设 $T_s$ 是 $H_s$ 的极大\Term{center}{中心}\Term{torus}{环面}（$s$ 是 $S$ 的\Term{closed-point}{闭点}），$M_\ell$ 是 $H$ 中\Term{lifting}{提升} ${}_\ell T_s$ 且满足 $H=\operatorname{Centr}_G M_\ell$ 的\Term{multiplicative-type}{乘法型}\Term{subgroup}{子群}。令 $T'_s$ 为 $T_s$ 在 $G'_s$ 中的像。由于 $G'$ 在 $S$ 上\Term{separated}{分离}（第 VI$_B$ 讲 5.2），$M_\ell$ 经 $u$ 的像是 $G'$ 的一个\Term{multiplicative-type}{乘法型}\Term{subgroup}{子群} $M'_\ell$（第 IX 讲 6.8）。置 $H'=\operatorname{Centr}_{G'}M'_\ell$，它是 $G'$ 的\Term{smooth}{光滑}\Term{group-subprescheme}{群子预概形}。对每个等于 $q$ 的幂的整数 $\ell'$，都存在 $\ell$ 使 $(M'_\ell)_s$ 包含 ${}_{\ell'}T'_s$，所以可以把 $\ell$ 取得充分大，使
\[H'_s=\operatorname{Centr}_{G'_s}T'_s=u_s(H_s),\]
最后一个等式来自 7.2。$u$ 在 $H$ 上的\Term{restriction}{限制} $v$ 显然经 $H'$ \Term{resolution}{分解}。证明 $v:H\to H'$ 是\Term{flat}{平坦}\Term{morphism}{态射}。由于 $H,H'$ 在 $S$ 上\Term{flat}{平坦}，且 $H_s\to H'_s$ \Term{flat}{平坦}，$v$ 在 $H_s$ 的某个\Term{neighborhood}{邻域}上\Term{flat}{平坦}（EGA IV 11.3.10 和 11.3.1）。故 $v$ 在 $H$ 的一个\Term{open-subgroup}{开子群}上\Term{flat}{平坦}（第 VI$_B$ 讲 2.2）；由于 $H$ 的\Term{fiber}{纤维}\Term{connected}{连通}，$v$ \Term{flat}{平坦}。于是 $H$ 的\Term{set}{集合}意义的像是 $H'$ 的\Term{open-subset}{开集}（必等于 $(H')^0$），赋予其\Term{induced-structure}{诱导结构}后，它表示（fpqc \Term{topology}{拓扑}下的）像\Term{sheaf}{层} $u(H)$。$u(H)\in\mathcal{CT}_{G'}(S)$ 随即由 7.2 得到。

再设 $H$ 是 $G$ 中\Term{torus}{环面} $T$ 的\Term{centralizer}{中心化子}，$T'$ 是 $T$ 经 $u$ 的像；它是 $G'$ 的\Term{subtorus}{子环面}（第 IX 讲 6.8）。$H$ 经 $u$ 的像包含于 $\operatorname{Centr}_{G'}(T')$，逐\Term{fiber}{纤维}与后者一致（7.2），且在 $S$ 上\Term{smooth}{光滑}，故 $u(H)=\operatorname{Centr}_{G'}(T')$。

\textbf{vi) 的证明。} 已知 $\mathcal{CT}_G$ 和 $\mathcal{CT}_{G'}$ 在 $S$ 上\Term{smooth}{光滑}（由 i)），为证明 $\widetilde u$ 是\Term{faithfully-flat}{忠实平坦}的 $S$-\Term{morphism}{态射}，只须在几何\Term{fiber}{纤维}上检验。因而归结到 $S$ 是\Term{algebraically-closed-field}{代数闭域} $k$ 的谱。设 $H'\in\mathcal{CT}_{G'}(k)$，$T'$ 是其极大\Term{center}{中心}\Term{torus}{环面}，$T$ 是 $G$ 中像为 $T'$ 的\Term{subtorus}{子环面}，$H=\operatorname{Centr}_G(T)$，从而 $H'=u(H)$（7.2）；再设 $N=\operatorname{Norm}_G(H)$，$N'=\operatorname{Norm}_{G'}(H')$。在 iv) 中已证明 $G/N$（相应地 $G'/N'$）\Term{canonical}{典范}地等同于 $\mathcal{CT}_G$ 中 $H$ 的\Term{open-neighborhood}{开邻域}
\SourcePage{308}
（相应地 $\mathcal{CT}_{G'}$ 中 $H'$ 的\Term{open-neighborhood}{开邻域}）。在这些等同下，$\widetilde u$ 在 $G/N$ 上的\Term{restriction}{限制}就是由 $u$ 取商得到的自然\Term{morphism}{态射}
\[w:G/N\longrightarrow G'/N'.\]
而 $w$ 是 $G$ 的\Term{homogeneous-space}{齐性空间}之间的\Term{epimorphism}{满态射}，所以\Term{faithfully-flat}{忠实平坦}。因此 $\widetilde u$ \Term{flat}{平坦}，且其像 $\widetilde u(\mathcal{CT}_G)$ 是 $\mathcal{CT}_{G'}$ 的\Term{open-subset}{开集}，包含 $\mathcal{CT}_{G'}(k)$ 的每个点。由于 $\mathcal{CT}_{G'}$ 在 $k$ 上\Term{finite-type}{有限型}，可知 $\widetilde u$ \Term{surjectivity}{满射}，故\Term{faithfully-flat}{忠实平坦}。

vi) 中关于 $\operatorname{Ker}u$ 为\Term{central}{中心的}情形的补充断言，来自第 XII 讲 7.12。

下面的\Term{theorem}{定理}推广第 XII 讲的\Term{theorem}{定理} 7.1。
\begin{SGAresult}{\Term{theorem}{定理}}{7.3}
设 $S$ 是\Term{prescheme}{预概形}，$G$ 是在 $S$ 上\Term{finite-presentation}{有限呈示}、\Term{smooth}{光滑}且\Term{fiber}{纤维}\Term{connected}{连通}的 $S$-\Term{group-prescheme}{群预概形}，考虑如下 $S$-\Term{functor}{函子} $\mathcal C:(\mathrm{Sch}/S)^\circ\to\mathrm{Ens}$：
\[\mathcal C(S')=G_{S'}\text{ 的所有 Cartan 子群所成的集合}.\]
\begin{enumerate}
\item[i)] 下列条件\Term{equivalence}{等价}：
\begin{enumerate}
\item[a)] \Term{functor}{函子} $\mathcal C$ \Term{representable}{可表}。
\item[b)] \Term{functor}{函子} $\mathcal C$ 可由在 $S$ 上\Term{smooth}{光滑}、\Term{quasi-projective}{拟射影}、\Term{finite-presentation}{有限呈示}且\Term{fiber}{纤维}\Term{affine}{仿射}的 $S$-\Term{prescheme}{预概形}表示。
\item[c)] 群 $G$ 对\Term{etale-topology}{平展拓扑}\Term{local-scheme}{局部}有一个\Term{cartan-subgroup}{Cartan 子群}。
\item[d)] 群 $G$ 对\Term{faithfully-flat-topology}{忠实平坦拓扑}\Term{local-scheme}{局部}有一个\Term{cartan-subgroup}{Cartan 子群}。
\item[e)] $G$ 的\Term{fiber}{纤维}的\Term{nilpotent-rank}{幂零秩}是 $S$ 上的\Term{locally-constant}{局部常值}函数。
\end{enumerate}
\item[ii)] 若上述条件成立，则 $G$ 的任意两个\Term{cartan-subgroup}{Cartan 子群}对\Term{etale-topology}{平展拓扑}\Term{local-scheme}{局部}\Term{conjugacy}{共轭}。$G$ 的\Term{fiber}{纤维}中所有\Term{regular-element}{正则点}（第 XIII 讲 2.7）所成的\Term{set}{集合}是一个在 $S$ 上\Term{finite-presentation}{有限呈示}的\Term{open-subset}{开集} $G_{\mathrm{r\text{é}g}}$，而 $G_{\mathrm{r\text{é}g}}$ 在 $S$ 上的每个\Term{section}{截面}都包含于 $G$ 的唯一一个\Term{cartan-subgroup}{Cartan 子群}。
\item[iii)] 设 $G'$ 是在 $S$ 上\Term{finite-presentation}{有限呈示}的 $S$-\Term{group-prescheme}{群预概形}，$u:G\to G'$ 是\Term{faithfully-flat}{忠实平坦}的 $S$-\Term{group-morphism}{群态射}，因而 $G'$ 在 $S$ 上\Term{smooth}{光滑}且\Term{fiber}{纤维}\Term{connected}{连通}。若 $C$ 是 $G$ 的\Term{cartan-subgroup}{Cartan 子群}，则 $u(C)=C'$ 可由 $G'$ 的一个\Term{cartan-subgroup}{Cartan 子群}表示，且 $C\to C'$ \Term{faithfully-flat}{忠实平坦}。
\item[iv)] 在 i) 与 iii) 的条件下，\Term{morphism}{态射}
\[\widetilde u:\mathcal C_G\longrightarrow\mathcal C_{G'},\qquad C\longmapsto u(C)=C'\]
\Term{faithfully-flat}{忠实平坦}。若再有 $\operatorname{Ker}(u)$ 是\Term{central}{中心的}，则 $\widetilde u$ 是\Term{isomorphism}{同构}。
\item[v)] 关于\Term{transporter}{输运子}及其与\Term{maximal-torus}{极大环面}的关系等补充结论，可参阅 7.1 及第 XII 讲 7.1。
\end{enumerate}
\end{SGAresult}
\SourcePage{309}
\begin{proof}
i) 证明 $\mathrm{b)}\Rightarrow\mathrm{c)}\Rightarrow\mathrm{d)}\Rightarrow\mathrm{e)}\Rightarrow\mathrm{b)}\Rightarrow\mathrm{a)}\Rightarrow\mathrm{d)}$。

$\mathrm{b)}\Rightarrow\mathrm{c)}$。设 $s\in S$。由于 $\mathcal C_s$ 在 $\kappa(s)$ 上\Term{smooth}{光滑}，$\mathcal C_s$ 中存在\Term{residue-field}{剩余域}为 $\kappa(s)$ 的\Term{finite-separable-extension}{有限可分扩张}的点。应用第 XI 讲 1.10，可知存在 $s$ 的\Term{open-neighborhood}{开邻域} $U$ 及\Term{surjectivity}{满射}\Term{etale}{平展}\Term{morphism}{态射} $U'\to U$，使 $G_{U'}$ 有一个\Term{cartan-subgroup}{Cartan 子群}。

$\mathrm{c)}\Rightarrow\mathrm{d)}$ 显然。

$\mathrm{d)}\Rightarrow\mathrm{e)}$。设 $s\in S$。由假设，存在在 $S$ 上\Term{flat}{平坦}、像包含 $s$ 的 $S$-\Term{prescheme}{预概形} $S'$，使 $G_{S'}$ 有一个\Term{cartan-subgroup}{Cartan 子群}。设 $s'\in S'$ 位于 $s$ 之上。则 $G_{S'}$ 的\Term{fiber}{纤维}的\Term{nilpotent-rank}{幂零秩}在 $\operatorname{Spec}\mathcal O_{S',s'}$ 上常值，因此 $G$ 的\Term{fiber}{纤维}的\Term{nilpotent-rank}{幂零秩}在其像 $\operatorname{Spec}\mathcal O_{S,s}$ 上常值（EGA IV 2.3.4 ii)）。设此值为 $r$。由 6.3 bis，$S$ 中满足 $G_x$ 的\Term{nilpotent-rank}{幂零秩}等于 $r$ 的点 $x$ 所成的\Term{set}{集合} $E_r$ 是 $S$ 的\Term{ind-constructible}{归纳可构造}部分，故包含 $s$ 的一个\Term{neighborhood}{邻域}（EGA IV 1.10.1）。

$\mathrm{e)}\Rightarrow\mathrm{b)}$。断言对 $S$ 是\Term{local-scheme}{局部}的，故可假定 $G$ 的\Term{fiber}{纤维}的\Term{nilpotent-rank}{幂零秩}恒等于 $r$。对于每个 $S$-\Term{prescheme}{预概形} $S'$，$G_{S'}$ 的\Term{cartan-subgroup}{Cartan 子群}恰好就是其中在 $S'$ 上\Term{smooth}{光滑}、\Term{relative-dimension}{相对维数}为 $r$，且几何\Term{fiber}{纤维}为某个\Term{torus}{环面}的\Term{centralizer}{中心化子}的\Term{group-subprescheme}{群子预概形}。由 7.1，$\mathcal{CT}$ \Term{representable}{可表}，所以 $\mathcal C$ 可由 $\mathcal{CT}$ 中表示\Term{relative-dimension}{相对维数}为 $r$ 的群所成的子\Term{functor}{函子}的既开又闭\Term{subprescheme}{子预概形}表示。b) 中其余断言均包含于 7.1 i)，只有 $\mathcal C$ 的\Term{fiber}{纤维}\Term{affine}{仿射}这一点来自第 XII 讲 7.1 d)。

$\mathrm{b)}\Rightarrow\mathrm{a)}$ 显然。证明 $\mathrm{a)}\Rightarrow\mathrm{d)}$。由 6.2 iii)，\Term{functor}{函子} $\mathcal C$ 与环的\Term{filtered}{滤过}\Term{inductive-limit}{归纳极限}交换，故若 $\mathcal C$ \Term{representable}{可表}，它必由一个在 $S$ 上\Term{locally-finite-presentation}{局部有限呈示}的 $S$-\Term{prescheme}{预概形}表示（EGA IV 8.14.2）。为证明 $\mathcal C$ 在 $S$ 上\Term{smooth}{光滑}，归结到证明：若 $S$ \Term{affine}{仿射}，$S_0$ 是由\Term{nilpotent-ideal}{幂零理想} $J$ 定义的闭\Term{subscheme}{子概形}，则 $G_0=G\times_S S_0$ 的每个\Term{cartan-subgroup}{Cartan 子群} $C_0$ 都\Term{liftable}{可提升}为 $G$ 的\Term{cartan-subgroup}{Cartan 子群} $C$。$C_0$ 的存在蕴含 $S_0$ 满足条件 e)，因 $S$ 与 $S_0$ 底\Term{space}{空间}相同，$S$ 也满足此条件，再由 $\mathrm{d)}\Rightarrow\mathrm{b)}$ 得到结论。由于 $\mathcal C$ 在 $S$ 上\Term{smooth}{光滑}且 $\mathcal C\to S$ \Term{surjectivity}{满射}，可见 $\mathrm{a)}\Rightarrow\mathrm{d)}$。

ii) 设 $C,C'$ 是 $G$ 的两个\Term{cartan-subgroup}{Cartan 子群}。则 $\operatorname{Transp}_G(C,C')$ 可由在 $S$ 上\Term{smooth}{光滑}（7.1 iii)）且\Term{fiber}{纤维}\Term{nonempty}{非空}（参见第 XII 讲 6.6 a)、c)）的\Term{prescheme}{预概形}表示。由 Hensel \Term{lemma}{引理}（第 XI 讲 1.10），$C,C'$ 对\Term{etale-topology}{平展拓扑}\Term{local-scheme}{局部}\Term{conjugacy}{共轭}。结合 i)，ii) 的其余断言来自第 XIII 讲 3.1 和 3.2。

iii) 设 $C$ 是 $G$ 的\Term{cartan-subgroup}{Cartan 子群}。由 7.1 v)，$u(C)$ 可由 $G'$ 的一个\Term{smooth}{光滑}\Term{subgroup}{子群} $C'$ 表示。由于 $C'$ 的\Term{fiber}{纤维}是 $G'$ 的\Term{fiber}{纤维}的\Term{cartan-subgroup}{Cartan 子群}（第 XII 讲 6.6 d)），$C'$ 是 $G'$ 的\Term{cartan-subgroup}{Cartan 子群}。

iv) 如 7.1 vi) 那样证明\Term{morphism}{态射} $\widetilde u$ \Term{faithfully-flat}{忠实平坦}。再设 $\operatorname{Ker}u$ 是\Term{central}{中心的}，$C'$ 是 $G'$ 的\Term{cartan-subgroup}{Cartan 子群}，则 $u^{-1}(C')=C$ \Term{smooth}{光滑}且\Term{fiber}{纤维}\Term{connected}{连通}（7.1 vi)），其\Term{fiber}{纤维}是\Term{cartan-subgroup}{Cartan 子群}（第 XII 讲 6.6 f)），
\SourcePage{310}
故 $C$ 是 $G$ 的\Term{cartan-subgroup}{Cartan 子群}。再结合 7.1 vi)，便完成 iv) 的证明。
\end{proof}

\SGAsection{8}{光滑群的子环面函子的可表性判据}
\textbf{8.0。} 本节中，若 $S$ 是\Term{prescheme}{预概形}，$G$ 是 $S$-\Term{group-prescheme}{群预概形}，则以 $\mathcal T_G$（无歧义时简记为 $\mathcal T$）表示如下 $S$-\Term{functor}{函子} $(\mathrm{Sch}/S)^\circ\to\mathrm{Ens}$：对每个 $S$-\Term{prescheme}{预概形} $S'$，
\[\mathcal T(S')=G_{S'}\text{ 的所有子环面所成的集合}.\]
类似地，以 $\mathcal{TC}$ 表示 $G$ 的\Term{center}{中心}\Term{subtorus}{子环面}\Term{functor}{函子}。
\begin{SGAresult}{命题}{8.1}
设 $S$ 是\Term{locally-noetherian}{局部 noether}\Term{prescheme}{预概形}，$G$ 是\Term{finite-type}{有限型} $S$-\Term{group-prescheme}{群预概形}。则下列条件\Term{equivalence}{等价}：
\begin{enumerate}
\item[i)] \Term{functor}{函子} $\mathcal T_G$“与 artin \Term{local-ring}{局部环}的进\Term{limit}{极限}交换”，即对任意形如 $S'=\operatorname{Spec}(A)$ 的 $S$-\Term{prescheme}{预概形}，其中 $A$ 为对其\Term{maximal-ideal}{极大理想} $\mathfrak m$ 所定义的\Term{topology}{拓扑}\Term{complete}{完备}的 noether \Term{local-ring}{局部环}，\Term{canonical}{典范}映射
\[\mathcal T(S')\longrightarrow\varprojlim_n\mathcal T(S'_n),\qquad S'_n=\operatorname{Spec}(A/\mathfrak m^n),\]
是\Term{bijection}{双射}。
\item[ii)] 同 i)，但只考虑 $A$ 为\Term{residue-field}{剩余域}\Term{algebraically-closed}{代数闭}的\Term{complete}{完备}\Term{discrete-valuation-ring}{离散赋值环}的情形。
\item[iii)] 同 ii)，但只考虑 $\mathcal T$ 中对应于 $G$ 的\Term{relative-dimension}{相对维数}为 $1$ 的\Term{subtorus}{子环面}的子\Term{functor}{函子} $\mathcal T^{(1)}$。
\item[i bis)] 对任意如 i) 中的 $S$-\Term{prescheme}{预概形} $S'$ 及任意 $S'$-\Term{torus}{环面} $T$，\Term{canonical}{典范}映射
\[\operatorname{Hom}_{S'\text{-}\mathrm{gr}}(T,G_{S'})\longrightarrow\varprojlim_n\operatorname{Hom}_{S'_n\text{-}\mathrm{gr}}(T_{S'_n},G_{S'_n})\]
是\Term{bijection}{双射}。
\item[ii bis)] 同 i bis)，但只考虑 $A$ 为\Term{residue-field}{剩余域}\Term{algebraically-closed}{代数闭}的\Term{complete}{完备}\Term{discrete-valuation-ring}{离散赋值环}的情形。
\item[iii bis)] 同 ii bis)，但只考虑 $T$ 为\Term{multiplicative-group}{乘法群} $\mathbb G_m$ 的情形。
\end{enumerate}
\end{SGAresult}
\begin{SGAresult}{注记}{8.2}
若只考虑 $G$ 的\Term{center}{中心}\Term{subtorus}{子环面}以及\Term{torus}{环面}到 $G$ 的\Term{central-homomorphism}{中心同态}，也有类似命题。
\end{SGAresult}
证明中将使用以下\Term{lemma}{引理}。
\begin{SGAresult}{\Term{lemma}{引理}}{8.3}
设 $S$ 是\Term{prescheme}{预概形}，$G$ 是 $S$-\Term{group-prescheme}{群预概形}，$T$ 是 $S$-\Term{torus}{环面}，$u:T\to G$ 是 $S$-\Term{group-morphism}{群态射}。再假定 $G$ 在 $S$ 上\Term{finite-presentation}{有限呈示}，或者 $S$ \Term{locally-noetherian}{局部 noether}且 $G$ \Term{locally-finite-type}{局部有限型}。则：
\begin{enumerate}
\item[a)] $\operatorname{Ker}u$ 是 $T$ 的\Term{multiplicative-type}{乘法型}\Term{subgroup}{子群}。
\item[b)] 商 $T'=T/\operatorname{Ker}u$ 是\Term{torus}{环面}。
\item[c)] 由 $u$ 取商得到的\Term{canonical}{典范}\Term{monomorphism}{单态射} $T'\to G$ 是\Term{immersion}{浸入}。
\end{enumerate}
\end{SGAresult}
\SourcePage{311}
当 $G$ 在 $S$ 上\Term{separated}{分离}时，此\Term{lemma}{引理}是第 IX 讲 6.8 的推论。一般情形按通常方法归结到 $S$ 是 noether 的情形。先证明 $K=\operatorname{Ker}u$ \Term{flat}{平坦}。可以假定 $S$ 是 artin \Term{local-ring}{局部环}的谱（EGA $0_{\mathrm{III}}$ 10.2.6），此时 $G$ \Term{separated}{分离}（第 VI$_B$ 讲 5.2），故 $K$ 为\Term{multiplicative-type}{乘法型}（第 IX 讲 6.8），从而在 $S$ 上\Term{flat}{平坦}。再证明 $\operatorname{Ker}u$ 在 $T$ 中闭；这归结到 $S$ 为\Term{discrete-valuation-ring}{离散赋值环}的谱的情形（EGA II 7.2.1）。由于 $T$ \Term{flat}{平坦}且\Term{fiber}{纤维}\Term{connected}{连通}，$u$ 经 $G$ 的泛\Term{fiber}{纤维}在 $G$ 中的\Term{schematic-closure}{概形闭包}的\Term{identity-component}{单位连通分支}\Term{resolution}{分解}。因此可以假定 $G$ \Term{flat}{平坦}且\Term{fiber}{纤维}\Term{connected}{连通}；此时 $G$ \Term{separated}{分离}（第 VI$_B$ 讲 5.2），故 $K$ 闭。由第 X 讲 4.8 b)，$K$ 是 $T$ 的\Term{multiplicative-type}{乘法型}\Term{subgroup}{子群}。于是商 $T'=T/K$ \Term{representable}{可表}，且 $T'$ 为\Term{multiplicative-type}{乘法型}群（第 IX 讲 2.7），其\Term{fiber}{纤维}为\Term{torus}{环面}，所以它是\Term{torus}{环面}。\Term{monomorphism}{单态射} $T'\to G$ 是\Term{immersion}{浸入}则来自第 VIII 讲 7.9。

\textbf{8.1 的证明。}

$\mathrm{i)}\Rightarrow\mathrm{i\ bis)}$。置 $T_n=T_{S'_n}$，$G_n=G_{S'_n}$，设
\[(u_n)_{n\in\mathbb N'}\in\varprojlim_n\operatorname{Hom}_{S'_n\text{-}\mathrm{gr}}(T_n,G_n).\]
对每个整数 $n$，$u_n(T_n)$ 都是 $G_n$ 的\Term{subtorus}{子环面} $T'_n$（\Term{lemma}{引理} 8.3）。由假设，存在唯一的 $G$ 的\Term{subtorus}{子环面} $T'$，对每个 $n$ 都\Term{lifting}{提升} $T'_n$。由于 $T'$ 的\Term{fiber}{纤维}\Term{affine}{仿射}，由 4.4 得到结论。

$\mathrm{i\ bis)}\Rightarrow\mathrm{i)}$。设 $(T_n)_{n\in\mathbb N}\in\varprojlim_n\mathcal T(S_n)$。由第 X 讲 4.6，存在一个 $S'$-\Term{torus}{环面} $T'$ 及 $S'_0$-\Term{isomorphism}{同构}
\[u_0:T'_0\xrightarrow{\sim}T_0.\]
由于 $T_n$ 在 $S'_n$ 上\Term{smooth}{光滑}，$u_0$ \Term{liftable}{可提升}为一个 $S'_n$-\Term{morphism}{态射}
\[u_n:T'_n\longrightarrow T_n\qquad\text{（第 IX 讲 3.6）},\]
且可令族 $(u_n)_{n\in\mathbb N}$ 相容，从而来自一个\Term{morphism}{态射} $u:T'\to G$。于是 $T'$ 经 $u$ 的像是 $G$ 的一个\Term{subtorus}{子环面}（\Term{lemma}{引理} 8.3），对每个 $n$ 都\Term{lifting}{提升} $T_n$。

蕴含 $\mathrm{i)}\Rightarrow\mathrm{ii)}\Rightarrow\mathrm{iii)}$ 及 $\mathrm{i\ bis)}\Rightarrow\mathrm{ii\ bis)}\Rightarrow\mathrm{iii\ bis)}$ 都显然。$\mathrm{iii)}\Rightarrow\mathrm{iii\ bis)}$ 的证明与 $\mathrm{i)}\Rightarrow\mathrm{i\ bis)}$ 相同。因此只须证明 $\mathrm{iii\ bis)}\Rightarrow\mathrm{ii\ bis)}\Rightarrow\mathrm{i\ bis)}$。

$\mathrm{ii\ bis)}\Rightarrow\mathrm{i\ bis)}$。采用 4.3 的术语，i bis) 成立的充要条件是 $\varprojlim_n\operatorname{Hom}_{S_n\text{-}\mathrm{gr}}(T_n,G_n)$ 的每个元素 $(u_n)_{n\in\mathbb N}$ 都“容许”（admissible）。对 $S'$ 中异于\Term{closed-point}{闭点} $s$ 的每个点 $t$，存在一个 $S$-\Term{scheme}{概形} $S''$，它是\Term{residue-field}{剩余域}\Term{algebraically-closed}{代数闭}的\Term{complete}{完备}\Term{discrete-valuation-ring}{离散赋值环}的谱，其\Term{generic-point}{泛点}映到 $t$，\Term{closed-point}{闭点}映到 $s$（EGA II 7.1.9）。这立即给出一个判定\Term{morphism}{态射}族容许的\Term{valuative-criterion}{赋值判据}，也就是说 $\mathrm{ii\ bis)}\Rightarrow\mathrm{i\ bis)}$。

$\mathrm{iii\ bis)}\Rightarrow\mathrm{ii\ bis)}$。设 $T$ 是 $S$-\Term{torus}{环面}，其中 $S$ 是\Term{residue-field}{剩余域}\Term{algebraically-closed}{代数闭}的\Term{complete}{完备}\Term{discrete-valuation-ring}{离散赋值环}的谱。则\Term{torus}{环面} $T$ 平凡（第 X 讲 4.6），即对某个适当整数 $r$ \Term{isomorphism}{同构}于 $(\mathbb G_{m,S})^r$。设
\[(u_n)_{n\in\mathbb N}\in\varprojlim_n\operatorname{Hom}_{S_n\text{-}\mathrm{gr}}((\mathbb G_m)^r_{S_n},G_n).\]
由假设，$u_n$ 在 $(\mathbb G_m)^r_S$ 的每个因子上的\Term{restriction}{限制}
\SourcePage{312}
来自一个\Term{group-morphism}{群态射} $\mathbb G_{m,S}\to G$。这给出一个乘积\Term{morphism}{态射} $(\mathbb G_m)^r_S\to G^r$，再与由 $G$ 的\Term{composition-law}{合成律}定义的\Term{morphism}{态射} $G^r\to G$ 复合，得到一个\Term{morphism}{态射}
\[v:(\mathbb G_m)^r_S\longrightarrow G.\]
由于存在\Term{group-morphism}{群态射} $u_n$，显然 $u_n=v_n$。还须证明 $v$ 是\Term{group-morphism}{群态射}，即两个显然的\Term{morphism}{态射}
\[f,g:X=(\mathbb G_m)^r_S\times_S(\mathbb G_m)^r_S\longrightarrow G\]
相等。令 $Z$ 为 $f,g$ 的等值\Term{subscheme}{子概形}。由于 $(\mathbb G_m)^r_S$ 的\Term{fiber}{纤维}\Term{connected}{连通}，且 $v((\mathbb G_m)^r_S)$ 包含\Term{unit-section}{单位截面}，如 8.3 中那样，可见 $v$ 经 $G$ 的\Term{identity-component}{单位连通分支}\Term{resolution}{分解}，从而可假定 $G$ \Term{separated}{分离}（第 VI$_B$ 讲 5.2），故 $Z$ 闭。另一方面，$v_n$ 是\Term{group-morphism}{群态射}，所以对每个 $n$ 都有 $f_n=g_n$；故 $Z$ 包含 $X$ 的闭\Term{fiber}{纤维}的某个\Term{neighborhood}{邻域}（EGA I 10.9.4）。由于 $X$ 的\Term{fiber}{纤维}\Term{smooth}{光滑}且\Term{irreducible}{不可约}（第 IX 讲 4.6），$Z$ 在 $X$ 中\Term{schematically-dense}{概形稠密}，从而 $Z=X$，即 $f=g$。

\begin{SGAresult}{定义}{8.4}
设 $S$ 是\Term{locally-noetherian}{局部 noether}\Term{prescheme}{预概形}，$G$ 是\Term{finite-type}{有限型} $S$-\Term{group-prescheme}{群预概形}，$\mathscr S$ 为满足以下条件的 $S$-\Term{scheme}{概形} $S'$ 所成的\Term{set}{集合}：$S'$ 是\Term{residue-field}{剩余域}\Term{algebraically-closed}{代数闭}的\Term{complete}{完备}\Term{discrete-valuation-ring}{离散赋值环}的谱，\Term{closed-point}{闭点}为 $s$，\Term{generic-point}{泛点}为 $t$，且 $(G_t)^0_{\mathrm{r\text{é}d}}$ \Term{smooth}{光滑}并具有 Chevalley \Term{resolution}{分解}，即是由一个\Term{smooth}{光滑}\Term{connected}{连通}线性\Term{algebraic-subgroup}{代数子群} $L_t$ 对一个\Term{abelian-variety}{阿贝尔簇}所作的\Term{extension}{扩张}（此\Term{resolution}{分解}于是唯一）。若 $S'\in\mathscr S$，以 $G'$（相应地 $L'$）表示 $G_t$（相应地 $L_t$）在 $G_{S'}$ 中的\Term{schematic-closure}{概形闭包}。

在这些条件下，若对每个 $S'\in\mathscr S$，$L_s$ 与 $G'_s$ 有相同的\Term{reductive-rank}{约化秩}，则称
\begin{quote}
“$G$ 的阿贝尔部分不\Term{degeneration}{退化}为\Term{torus}{环面}部分”，或简称 $G$ 满足性质 AT。
\end{quote}
（直观地说，假定商 $A=G'/L$ \Term{representable}{可表}，此时 $A$ 是\Term{flat}{平坦}\Term{group-prescheme}{群预概形}，且 $(A_t)^0_{\mathrm{r\text{é}d}}$ 是\Term{abelian-variety}{阿贝尔簇}。条件“AT”就是说 $A_s$ 的\Term{reductive-rank}{约化秩}为零，故 $(A_s)^0_{\mathrm{r\text{é}d}}$ 是由\Term{unipotent-group}{幺幂群}对一个\Term{abelian-variety}{阿贝尔簇}所作的\Term{extension}{扩张}。）

类似地，再假定 $G$ 的\Term{fiber}{纤维}\Term{connected}{连通}，若对每个 $S'\in\mathscr S$，$G_t$ 的\Term{center}{中心} $Z_t$ 的\Term{schematic-closure}{概形闭包} $Z$ 满足 AT，则称 $G$ 满足性质 ATC。
\end{SGAresult}
以下命题说明这些技术性定义的依据。
\begin{SGAresult}{命题}{8.5}
设 $S$ 是\Term{locally-noetherian}{局部 noether}\Term{prescheme}{预概形}，$G$ 是\Term{finite-type}{有限型} $S$-\Term{group-prescheme}{群预概形}。则：
\begin{enumerate}
\item[a)] $G$ 的\Term{subtorus}{子环面}\Term{functor}{函子} $\mathcal T$“与 artin \Term{local-ring}{局部环}的进\Term{limit}{极限}交换”（参见 8.1）的充要条件是 $G$ 满足性质 AT（8.4）。
\SourcePage{313}
\item[b)] 若 $G$ 的\Term{fiber}{纤维}\Term{connected}{连通}，则 $G$ 的\Term{center}{中心}\Term{subtorus}{子环面}\Term{functor}{函子} $\mathcal{TC}$ 与 artin \Term{local-ring}{局部环}的进\Term{limit}{极限}交换的充要条件是 $G$ 满足性质 ATC（8.4）。
\end{enumerate}
\end{SGAresult}
需要以下技术性\Term{lemma}{引理}。
\begin{SGAresult}{\Term{lemma}{引理}}{8.6}
设 $S$ 是\Term{discrete-valuation-ring}{离散赋值环}的谱，$s$ 是 $S$ 的\Term{closed-point}{闭点}，$G$ 是\Term{flat}{平坦}\Term{finite-type}{有限型} $S$-\Term{group-prescheme}{群预概形}，$T_s$ 是 $G_s$ 的\Term{subtorus}{子环面}。则存在一个在 $S$ 上\Term{faithfully-flat}{忠实平坦}的 $S$-\Term{scheme}{概形} $S'$，它是\Term{discrete-valuation-ring}{离散赋值环}的谱，\Term{closed-point}{闭点}为 $s'$；并存在 $G_{S'}$ 的\Term{group-subscheme}{群子概形} $C$，它在 $S'$ 上\Term{flat}{平坦}、交换且\Term{fiber}{纤维}\Term{connected}{连通}，满足 $C_{s'}$ 包含 $T_s\times_s s'$。
\end{SGAresult}
将 $S$ 换成一个在其上\Term{faithfully-flat}{忠实平坦}的\Term{discrete-valuation-ring}{离散赋值环}的谱 $S'$ 后，可以假定 $T_s=(\mathbb G_m)^r_s$，且 $\kappa(s)$ 在其\Term{prime-field}{素域}上的\Term{transcendence-degree}{超越次数} $\geq r$（EGA $0_{\mathrm{III}}$ 10.3.1 和 EGA II 7.1.9）。于是存在 $\bar x\in T_s(s)$，使 $G_s$ 中任意“包含”$\bar x$ 的\Term{algebraic-subgroup}{代数子群}都包含 $T_s$（参见第 XIII 讲 2.1 的 $\mathrm{(ii)}\Rightarrow\mathrm{(vii)}$ 的证明）。由于 $G$ 在 $S$ 上\Term{flat}{平坦}，有一个经过 $\bar x$ 的\Term{quasi-section}{拟截面}（EGA IV 14.5.8）；所以再次用一个在 $S$ 上\Term{faithfully-flat}{忠实平坦}的\Term{discrete-valuation-ring}{离散赋值环}的谱代替 $S$ 后，可以假定存在 $G$ 在 $S$ 上\Term{lifting}{提升} $\bar x$ 的\Term{section}{截面} $x$。令 $C_t$ 为 $G_t$ 中由 $x_t$ 生成的\Term{commutative-algebra}{交换代数}\Term{subgroup}{子群}（第 VI$_B$ 讲 7），$C$ 为 $C_t$ 在 $G$ 中的\Term{schematic-closure}{概形闭包}。显然 $C_s$ 包含 $\bar x$，从而包含 $T_s$；因此 $C$ 的“\Term{identity-component}{单位连通分支}”便是所求的\Term{flat}{平坦}交换\Term{group-scheme}{群概形}。

\textbf{8.5 a) 的证明。}

先假定\Term{functor}{函子} $\mathcal T$ 与 artin 环的进\Term{limit}{极限}交换，证明 $G$ 满足性质 AT。设 $S'\in\mathscr S$，$T_s$ 是 $G'_s$ 的一个\Term{maximal-torus}{极大环面}。须证明 $T_s$ 包含于 $L_s$。$L$ 与 $G'$ 的形成显然与\Term{discrete-valuation-ring}{离散赋值环}的\Term{faithfully-flat}{忠实平坦}\Term{group-extension}{扩张} $S''\to S'$ 交换。因此由\Term{lemma}{引理} 8.6，改变 $S'$ 后，可以假定存在 $G'_{S'}$ 的一个\Term{flat}{平坦}交换\Term{group-subscheme}{群子概形} $C$，使 $C_s$ 包含 $T_s$。于是 $T_s$ 是 $C_s$ 的\Term{center}{中心}\Term{subtorus}{子环面}，因而可\Term{infinitesimal}{无穷小}地\Term{lifting}{提升}为\Term{center}{中心}\Term{subtorus}{子环面}（推论 2.3）。由对 $G$ 的假设，$T_s$ \Term{liftable}{可提升}为 $G$ 的\Term{subtorus}{子环面} $T$。显然 $T_t$ 包含于 $G_t$ 的线性部分 $L_t$，故 $T$ 包含于 $L$。

再假定 $G$ 满足性质 AT，证明条件 8.1 iii bis) 成立。设 $S$ 是\Term{residue-field}{剩余域}\Term{algebraically-closed}{代数闭}的\Term{complete}{完备}\Term{discrete-valuation-ring}{离散赋值环} $A$ 的谱，$\mathfrak m$ 为 $A$ 的\Term{maximal-ideal}{极大理想}，$S_n=\operatorname{Spec}(A/\mathfrak m^n)$，$u_n$（$n\in\mathbb N$）为一个相容的\Term{group-morphism}{群态射}系
\[u_n:(\mathbb G_{m,S})_n\longrightarrow G_n=G\times_S S_n.\]
取在 $S$ 上\Term{invertible}{可逆}的素数 $q$。当整数 $\ell$ 遍历 $q$ 的幂时，存在唯一的 $S$-\Term{group-morphism}{群态射}
\[{}_\ell u:{}_\ell\mathbb G_{m,S}\longrightarrow G\]
对每个 $n$ 都\Term{lifting}{提升} $u_n|_{{}_\ell(\mathbb G_{m,S})_n}$（命题 1.6 a)）。因此，若存在对每个 $n$ 都\Term{lifting}{提升} $u_n$ 的 $S$-\Term{group-morphism}{群态射} $u:\mathbb G_{m,S}\to T$，它在 ${}_\ell\mathbb G_{m,S}$ 上的\Term{restriction}{限制}便唯一确定。结合\Term{dense}{稠密}性\Term{theorem}{定理}（第 IX 讲 4.8 a)），这证明了
\SourcePage{314}
$u$ 的唯一性，并说明为证明 $u$ 的存在，可以对底作\Term{faithfully-flat}{忠实平坦}\Term{group-extension}{扩张}。而 $(G_t)^0_{\mathrm{r\text{é}d}}$ 具有 Chevalley \Term{resolution}{分解}，且此\Term{resolution}{分解}已在 $A$ 的\Term{field-of-fractions}{分式域} $K$ 的某个\Term{finite-extension}{有限扩张} $L$ 上定义。因此用 $S$ 在 $L$ 中的\Term{normalization}{正规化}代替 $S$ 后，可假定 $S\in\mathscr S$。\Term{morphism}{态射} ${}_\ell u_t$ 经 $(G_t)_{\mathrm{r\text{é}d}}$ \Term{resolution}{分解}，故 ${}_\ell u$ 经 $(G_t)_{\mathrm{r\text{é}d}}$ 在 $G'$ 中的\Term{schematic-closure}{概形闭包} $G''$ \Term{resolution}{分解}。再次由\Term{dense}{稠密}性\Term{theorem}{定理}，可知对每个 $n$，$u_n$ 经 $G''_n$ \Term{resolution}{分解}。由于 $G$ 具有性质 AT，$G'_s$ 的每个\Term{subtorus}{子环面}，从而 $G''_s$ 的每个\Term{subtorus}{子环面}都包含于 $L_s$，故 $u_s$ 经 $L_s$ \Term{resolution}{分解}。

断言对每个 $n$，$u_n$ 都经 $L_n$ \Term{resolution}{分解}。事实上，由于 $L_t$ 在 $(G'')_t$ 中\Term{normal}{正规}，$L$ 在 $G''$ 中\Term{normal}{正规}；又由于 $L$ 在 $S$ 上\Term{flat}{平坦}，对每个整数 $n$，商 $H_n=G''_n/L_n$ \Term{representable}{可表}（第 VI$_A$ 讲 \S4）。$(\mathbb G_{m,S})_n$ 在 $H_n$ 中的像是 $H_n$ 的\Term{subtorus}{子环面}（第 IX 讲 6.8），其闭\Term{fiber}{纤维}为零，故此像为零，因而 $u_n$ 经 $L_n$ \Term{resolution}{分解}。而 $L_t$ \Term{affine}{仿射}，所以可知族 $u_n$ 容许（命题 4.3），证明完毕。

结合 8.2 及第 IX 讲 5.6 a)，8.5 b) 的证明与 8.5 a) 完全类似。

\begin{SGAresult}{命题}{8.7}
设 $S$ 是\Term{locally-noetherian}{局部 noether}\Term{prescheme}{预概形}，$G$ 是\Term{finite-type}{有限型} $S$-\Term{group-prescheme}{群预概形}。则：
\begin{enumerate}
\item[i)] 若 $G$ 的\Term{fiber}{纤维}\Term{connected}{连通}且满足 AT，则它满足 ATC。
\item[ii)] 若 $G$ 满足 AT，则 $G$ 的每个\Term{group-subprescheme}{群子预概形}都满足 AT。
\item[iii)] 若 $G$ 在 $S$ 上\Term{flat}{平坦}，且 $G$ 的\Term{fiber}{纤维}的\Term{abelian-rank}{阿贝尔秩}是 $S$ 上的\Term{locally-constant}{局部常值}函数，则 $G$ 满足 AT。
\end{enumerate}
\end{SGAresult}
i) 立即来自 8.5 及第 IX 讲 5.6 a)。

ii) 设 $H$ 是 $G$ 的\Term{group-subprescheme}{群子预概形}。为证明 $H$ 满足性质 AT，可以假定 $S$ 是\Term{complete}{完备}\Term{discrete-valuation-ring}{离散赋值环}的谱，而须证明任意相容的 $S_n$-\Term{group-morphism}{群态射}族
\[u_n:(\mathbb G_{m,S})_n\longrightarrow H_n\]
均容许（8.5 及 8.1 iii bis)）。由假设，$G$ 满足性质 AT，故存在对每个 $n$ 都\Term{lifting}{提升} $u_n$ 的 $S$-\Term{group-morphism}{群态射}
\[u:(\mathbb G_m)_S\longrightarrow G.\]
如 8.5 的证明那样，可知 $u|_{{}_\ell\mathbb G_{m,S}}$ 经 $H$ \Term{resolution}{分解}，其中 $\ell$ 遍历某个在 $S$ 上\Term{invertible}{可逆}的素数 $q$ 的幂。由\Term{dense}{稠密}性可知，在泛\Term{fiber}{纤维}上 $u_t$ 经 $H_t$ \Term{resolution}{分解}。由于 $(\mathbb G_m)_S$ \Term{reduced}{约化}，这就说明 $u$ 经 $H$ \Term{resolution}{分解}。

iii) 设 $S'\in\mathscr S$（采用 8.4 的记号）。群 $L$ 在 $S$ 上\Term{flat}{平坦}，其泛\Term{fiber}{纤维} $L_t$ \Term{affine}{仿射}；在这些条件下，可以证明 $L_s$ 必\Term{affine}{仿射}
\SourcePage{315}
（第 XVII 讲附录 III，2）。由于 $G,L$ \Term{flat}{平坦}，有 $G=G'$，且 $G,L$ 的\Term{fiber}{纤维}\Term{dimension}{维数}在 $S$ 上常值（第 VI$_B$ 讲 4），因而有不等式
\[\operatorname{rg}_{\mathrm{ab}}G_s\leq\dim G_s-\dim L_s=\dim G_t-\dim L_t=\operatorname{rg}_{\mathrm{ab}}G_t.\]
由假设，这些不等式均为等式。因此 $(G_s/L_s)^0_{\mathrm{r\text{é}d}}$ 是\Term{abelian-variety}{阿贝尔簇}，故\Term{reductive-rank}{约化秩}为零，从而 $L_s$ 与 $G_s$ 有相同的\Term{reductive-rank}{约化秩}。

现在可以证明本节的主要\Term{theorem}{定理}。
\begin{SGAresult}{\Term{theorem}{定理}}{8.8}
设 $G$ 是\Term{locally-noetherian}{局部 noether}\Term{prescheme}{预概形} $S$ 上的\Term{finite-type}{有限型}\Term{group-prescheme}{群预概形}。假定 $G$ 在 $S$ 上\Term{flat}{平坦}且\Term{fiber}{纤维}\Term{connected}{连通}。则：
\begin{enumerate}
\item[a)] $G$ 的\Term{center}{中心}\Term{subtorus}{子环面}\Term{functor}{函子} $\mathcal{TC}$ \Term{representable}{可表}的充要条件是 $G$ 具有性质 ATC（8.4）。而且此时 $\mathcal{TC}$ 可由在 $S$ 上\Term{etale}{平展}且\Term{separated}{分离}的 $S$-\Term{prescheme}{预概形}表示。
\item[b)] 在 a) 的条件下，对每个 $S$-\Term{torus}{环面} $T$，由 $T$ 到 $G$ 的\Term{central-homomorphism}{中心同态}所成的\Term{functor}{函子} $\operatorname{Hom}^{\mathrm{cent}}_{S\text{-}\mathrm{gr}}(T,G)$ 可由在 $S$ 上\Term{etale}{平展}且\Term{separated}{分离}的 $S$-\Term{prescheme}{预概形}表示。
\end{enumerate}
\end{SGAresult}
\begin{SGAresult}{\Term{theorem}{定理}}{8.9}
设 $S$ 是\Term{locally-noetherian}{局部 noether}\Term{prescheme}{预概形}，$G$ 是在 $S$ 上\Term{smooth}{光滑}的\Term{finite-type}{有限型} $S$-\Term{group-prescheme}{群预概形}。
\begin{enumerate}
\item[a)] $G$ 的\Term{subtorus}{子环面}\Term{functor}{函子} $\mathcal T$ \Term{representable}{可表}的充要条件是 $G$ 具有性质 AT（8.4）。而且此时 $\mathcal T$ 可由在 $S$ 上\Term{smooth}{光滑}且\Term{separated}{分离}的 $S$-\Term{prescheme}{预概形}表示；更精确地说，\Term{structural-morphism}{结构态射} $\mathcal T\to S$ 有一个\Term{canonical}{典范}\Term{resolution}{分解}
\[\mathcal T\xrightarrow{u}Y\xrightarrow{v}S,\]
其中 $v$ 是\Term{smooth}{光滑}且\Term{quasi-projective}{拟射影}（从而\Term{finite-type}{有限型}）的\Term{morphism}{态射}，$u$ 是\Term{etale}{平展}\Term{separated}{分离}\Term{morphism}{态射}。
\item[b)] 在 a) 的条件下，对每个 $S$-\Term{torus}{环面} $T$，从 $T$ 到 $G$ 的\Term{homomorphism}{同态}\Term{functor}{函子} $\operatorname{Hom}_{S\text{-}\mathrm{gr}}(T,G)$ 可由在 $S$ 上\Term{smooth}{光滑}且\Term{separated}{分离}的 $S$-\Term{prescheme}{预概形}表示。
\end{enumerate}
\end{SGAresult}

\textbf{8.8 a) 的证明。} 若\Term{functor}{函子} $\mathcal{TC}$ \Term{representable}{可表}，则它与 artin 环的进\Term{limit}{极限}交换，从而（8.2 及 8.5 b)）$G$ 具有性质 ATC。为证明逆命题，将使用以下结果；它将在 EGA VI 中证明，也见于 Murre 的报告：Bourbaki 讨论班，1965 年 5 月，第 294 号，\Term{theorem}{定理} 1、推论 2。
\begin{SGAresult}{\Term{lemma}{引理}}{8.10}
设 $S$ 是\Term{locally-noetherian}{局部 noether}\Term{prescheme}{预概形}，$F$ 是定义在 $\mathrm{Sch}/S$ 上、取值于\Term{set}{集合}\Term{category}{范畴}的\Term{contravariant}{反变}\Term{functor}{函子}。为使 $F$ 可由\Term{etale}{平展}\Term{separated}{分离}的 $S$-\Term{prescheme}{预概形}表示，（必要且）充分的条件是 $F$ 满足以下五条性质：
\begin{enumerate}
\item[i)] $F$ 是 fpqc \Term{topology}{拓扑}的\Term{sheaf}{层}（第 IV 讲）。
\SourcePage{316}
\item[ii)] $F$ 与环的\Term{inductive-limit}{归纳极限}交换（第 XI 讲 3.2）。
\item[iii)] $F$ 与 artin \Term{local-ring}{局部环}的进\Term{limit}{极限}交换（8.1 i)）。
\item[iv)] $F$ 满足“\Term{valuative-criterion-separatedness}{分离性赋值判据}”，即对每个为\Term{discrete-valuation-ring}{离散赋值环}的谱、\Term{generic-point}{泛点}为 $t$ 的 $S$-\Term{scheme}{概形} $S'$，\Term{canonical}{典范}映射
\[F(S')\longrightarrow F(t)\]
是\Term{injectivity}{单射}。
\item[v)] $F$ 是\Term{infinitesimally-etale}{无穷小平展}的（第 XI 讲 1.8）。
\end{enumerate}
\end{SGAresult}
证明 8.8 的\Term{functor}{函子} $\mathcal{TC}$ 满足 8.10 的五条条件。

i) 只要 $G$ 在 $S$ 上\Term{finite-presentation}{有限呈示}，\Term{functor}{函子} $\mathcal T$（相应地 $\mathcal{TC}$）就是 fpqc \Term{topology}{拓扑}的\Term{sheaf}{层}。事实上，此时\Term{torus}{环面}到 $G$ 的每个\Term{monomorphism}{单态射}都是\Term{immersion}{浸入}（第 VIII 讲 7.9），再用\Term{subprescheme}{子预概形}的 fpqc \Term{descent}{下降}便得结论。

ii) 推论 6.3 说明，当 $G$ 在 $S$ 上\Term{finite-presentation}{有限呈示}时，\Term{functor}{函子} $\mathcal T$ 与环的\Term{inductive-limit}{归纳极限}交换；立即可知 $\mathcal{TC}$ 也如此。

iii) 由 8.5，条件 iii) 恰好\Term{equivalence}{等价}于性质 ATC。

iv) 只须注意：若 $S$ 是\Term{discrete-valuation-ring}{离散赋值环}的谱，则 $G$ 中泛\Term{fiber}{纤维}相同的两个\Term{subtorus}{子环面}必相同；更精确地说，它们都等于其泛\Term{fiber}{纤维}在 $G$ 中的\Term{schematic-closure}{概形闭包}的\Term{identity-component}{单位连通分支}。

v) 由 2.3 得到，因为 $G$ 在 $S$ 上\Term{flat}{平坦}。

\textbf{8.8 b) 的证明。} 仿照第 XI 讲 4.2，可见只须证明乘积群 $T\times_S G$ 仍满足性质 ATC，而由定义立即可知。

\textbf{8.9 a) 的证明。} 用 $G$ 的\Term{identity-component}{单位连通分支}代替 $G$（第 VI$_B$ 讲 3.10），可以假定 $G$ 的\Term{fiber}{纤维}\Term{connected}{连通}。若 $T$ 是 $G$ 的\Term{subtorus}{子环面}，则其\Term{centralizer}{中心化子} $\operatorname{Centr}_G(T)$ \Term{representable}{可表}（第 XI 讲 6.11），且是 $G$ 中在 $S$ 上\Term{smooth}{光滑}（第 XI 讲 2.4）、\Term{fiber}{纤维}\Term{connected}{连通}（第 XII 讲 6.6 b)）的\Term{group-subprescheme}{群子预概形}，故属于 $\mathcal{CT}(S)$；这里 $\mathcal{CT}$ 是 7.1 i) 定义的\Term{functor}{函子}。显然映射 $T\mapsto\operatorname{Centr}_G(T)$ 定义了一个 $S$-\Term{morphism}{态射}
\[u:\mathcal T\longrightarrow\mathcal{CT}.\]
由于 $\mathcal{CT}$ 可由在 $S$ 上\Term{smooth}{光滑}且\Term{quasi-projective}{拟射影}的 $S$-\Term{prescheme}{预概形}表示（7.1 i)），只须证明\Term{morphism}{态射} $u$ 可由\Term{separated}{分离}\Term{etale}{平展}\Term{morphism}{态射}表示。

作适当的\Term{base-change}{基变换} $S'\to S$（其中 $S'=\mathcal{CT}$，故 $S'$ \Term{locally-noetherian}{局部 noether}）后，归结到以下问题：设 $S$ 是\Term{locally-noetherian}{局部 noether}\Term{prescheme}{预概形}，$G$ 是在 $S$ 上\Term{smooth}{光滑}且\Term{finite-type}{有限型}、\Term{fiber}{纤维}\Term{connected}{连通}的 $S$-\Term{group-prescheme}{群预概形}，
\SourcePage{317}
$H$ 是 $G$ 中在 $S$ 上\Term{smooth}{光滑}且\Term{fiber}{纤维}\Term{connected}{连通}的\Term{subgroup}{子群}。考虑 $\mathcal{TC}_H$ 的子\Term{functor}{函子} $X$：对每个 $S$-\Term{prescheme}{预概形} $S'$，$X(S')$ 是满足 $\operatorname{Centr}_{G_{S'}}(T)=H_{S'}$ 的 $H_{S'}$ 的\Term{center}{中心}\Term{subtorus}{子环面} $T$ 所成的\Term{set}{集合}。将证明 $X$ 可由在 $S$ 上\Term{etale}{平展}且\Term{separated}{分离}的 $S$-\Term{prescheme}{预概形}表示。

事实上，由假设 $G$ 满足性质 AT，故 $H$ 也满足（8.7 ii)）；由于 $H$ 的\Term{fiber}{纤维}\Term{connected}{连通}，它也满足性质 ATC（8.7 i)）。另一方面，$H$ 在 $S$ 上\Term{smooth}{光滑}，故\Term{flat}{平坦}。由 8.8 a)，$\mathcal{TC}_H$ 可由在 $S$ 上\Term{etale}{平展}且\Term{separated}{分离}的 $S$-\Term{prescheme}{预概形}表示。因此只须证明\Term{canonical}{典范}\Term{monomorphism}{单态射} $X\to\mathcal{TC}_H$ 可由\Term{open-immersion}{开浸入}表示。

置 $S'=\mathcal{TC}_H$，令 $K$ 为 $H_{S'}$ 的“泛”\Term{center}{中心}\Term{torus}{环面}在 $G_{S'}$ 中的\Term{centralizer}{中心化子}。群 $K$ 是在 $S'$ 上\Term{smooth}{光滑}、\Term{fiber}{纤维}\Term{connected}{连通}且包含 $H_{S'}$ 的\Term{group-scheme}{群概形}。由定义，
\[X=\prod_{K/S'}H_{S'}/K,\]
它确可由 $S'$ 中使 $H_{S'}$ 与 $K$ 有相同\Term{relative-dimension}{相对维数}的既开又闭\Term{subprescheme}{子预概形}表示。

8.9 b) 的证明与 8.8 b) 类似。

\begin{SGAresult}{推论}{8.11}
设 $S$ 是\Term{prescheme}{预概形}，$G$ 是在 $S$ 上\Term{smooth}{光滑}且\Term{finite-presentation}{有限呈示}的 $S$-\Term{group-prescheme}{群预概形}。若 $G$ 的\Term{fiber}{纤维}的\Term{abelian-rank}{阿贝尔秩}是 $S$ 上的\Term{locally-constant}{局部常值}函数（特别地，若 $G$ 的\Term{fiber}{纤维}\Term{affine}{仿射}），则 $G$ 的\Term{subtorus}{子环面}\Term{functor}{函子}可由在 $S$ 上\Term{smooth}{光滑}且\Term{separated}{分离}的 $S$-\Term{prescheme}{预概形}表示；对任意 $S$-\Term{torus}{环面} $T$，\Term{functor}{函子} $\operatorname{Hom}_{S\text{-}\mathrm{gr}}(T,G)$ 也如此。
\end{SGAresult}
事实上，断言对 $S$ 是\Term{local-scheme}{局部}的，故可以假定 $S$ \Term{affine}{仿射}且 $G$ 的\Term{fiber}{纤维}的\Term{abelian-rank}{阿贝尔秩}常值。于是可以找到一个\Term{affine}{仿射} noether \Term{scheme}{概形} $S_0$ 和一个在 $S_0$ 上\Term{smooth}{光滑}、\Term{finite-type}{有限型}的 $S_0$-\Term{group-prescheme}{群预概形} $G_0$，使 $G$ 与 $G_0\times_{S_0}S$ 为 $S$-\Term{isomorphism}{同构}（第 VI$_B$ 讲 \S10）。另一方面，$S$ 上\Term{finite-presentation}{有限呈示}的 $S$-\Term{group-prescheme}{群预概形}的\Term{fiber}{纤维}的\Term{abelian-rank}{阿贝尔秩}是\Term{ind-constructible}{归纳可构造}函数（6.3 bis）。由标准论证（EGA IV 8），在当前情形下，可以假定 $G_0$ 的\Term{fiber}{纤维}的\Term{abelian-rank}{阿贝尔秩}在 $S_0$ 上常值。于是 $G_0$ 具有性质 AT（8.7 iii)），从而（8.10）$G_0$ 的\Term{subtorus}{子环面}\Term{functor}{函子}可由在 $S_0$ 上\Term{smooth}{光滑}且\Term{separated}{分离}的 $S_0$-\Term{prescheme}{预概形}表示；这就得到 $G$ 的所述性质。对于\Term{functor}{函子} $\operatorname{Hom}_{S\text{-}\mathrm{gr}}(T,G)$，作类似处理。

\textbf{8.9 的推广。} \Term{smooth}{光滑}群 $G$ 的\Term{subtorus}{子环面}\Term{functor}{函子} $\mathcal T$ 未必\Term{representable}{可表}，因此下面给出 $\mathcal T$ 的子\Term{functor}{函子}\Term{representable}{可表}的充分条件。
\begin{SGAresult}{命题}{8.12}
设 $S$ 是\Term{locally-noetherian}{局部 noether}\Term{prescheme}{预概形}，$G$ 是在 $S$ 上\Term{smooth}{光滑}且\Term{fiber}{纤维}\Term{connected}{连通}的 $S$-\Term{group-prescheme}{群预概形}，$F$ 是 $G$ 的\Term{subtorus}{子环面}\Term{functor}{函子} $\mathcal T$ 的一个子 $S$-\Term{functor}{函子}，满足以下性质：
\begin{enumerate}
\item[i)] $F$ 是 fpqc \Term{topology}{拓扑}的\Term{sheaf}{层}（第 IV 讲）。
\item[ii)] $F$ 与环的\Term{inductive-limit}{归纳极限}交换（第 XI 讲 3.2）。
\SourcePage{318}
\item[iii)] $F$ 与 artin \Term{local-ring}{局部环}的进\Term{limit}{极限}交换（第 XI 讲 3.3）。
\item[iv)] $F$ 在 $S$ 上\Term{infinitesimally-smooth}{无穷小光滑}（第 XI 讲 1.8）。
\item[v)] $F$ 在 $G$ 的\Term{inner-automorphism}{内自同构}下\Term{stable}{稳定}，即对每个 $S$-\Term{prescheme}{预概形} $S'$，有
\[T\in F(S')\ \text{且}\ g\in G(S')\quad\Longrightarrow\quad\operatorname{int}(g)T\in F(S').\]
\end{enumerate}
则 $F$ 可由在 $S$ 上\Term{smooth}{光滑}且\Term{separated}{分离}的 $S$-\Term{prescheme}{预概形}表示。
\end{SGAresult}
\begin{SGAresult}{命题}{8.12 bis}
设 $S,G$ 如上，$T$ 是 $S$-\Term{torus}{环面}，$F$ 是 $\operatorname{Hom}_{S\text{-}\mathrm{gr}}(T,G)$ 的子\Term{functor}{函子}，满足以下性质：

上述 i)、ii)、iii)、iv)；

v) $F$ 在 $G$ 的\Term{inner-automorphism}{内自同构}下\Term{stable}{稳定}，即对每个 $S$-\Term{prescheme}{预概形} $S'$，有
\[u\in F(S')\ \text{且}\ g\in G(S')\quad\Longrightarrow\quad\operatorname{int}(g)u\in F(S').\]
则 $F$ 可由在 $S$ 上\Term{smooth}{光滑}且\Term{separated}{分离}的 $S$-\Term{prescheme}{预概形}表示。
\end{SGAresult}
\textbf{8.12 的证明。} （8.12 bis 的证明完全类似，留给读者。）

设 $u:F\to\mathcal{CT}$（7.1 i)）是将每个 $T\in F(S')$ 映到 $\operatorname{Centr}_{G_{S'}}(T)\in\mathcal{CT}(S')$ 的 $S$-\Term{morphism}{态射}。由于 $\mathcal{CT}$ 可由\Term{smooth}{光滑}\Term{quasi-projective}{拟射影}的 $S$-\Term{prescheme}{预概形}表示（7.1 i)），只须证明 $u$ 的\Term{representability}{可表性}。经\Term{base-change}{基变换} $\mathcal{CT}\to S$ 后，归结到以下问题：给定如上的 $S,G$ 及 $G$ 中\Term{smooth}{光滑}且\Term{fiber}{纤维}\Term{connected}{连通}的\Term{subgroup}{子群} $H$，须表示\Term{functor}{函子} $F_H$，其中 $F_H(S')$ 为满足
\[H_{S'}=\operatorname{Centr}_{G_{S'}}(T)\]
的 $T\in F(S')$ 所成的\Term{set}{集合}。

将证明 $F_H$ 在 $S$ 上\Term{etale}{平展}且\Term{separated}{分离}。为此只须核实 8.10 的五条条件。8.10 的条件 i)、ii)、iii) 容易由 8.12 i)、ii)、iii) 得到。已经指出 $\mathcal{TC}_H$ 是\Term{separated}{分离}且\Term{infinitesimally-etale}{无穷小平展}的\Term{functor}{函子}，所以作为 $\mathcal{TC}_H$ 的子\Term{functor}{函子}，$F_H$ \Term{separated}{分离}且\Term{infinitesimally-unramified}{无穷小非分歧}（第 XI 讲 1.8）。因此只须证明 $F_H$ \Term{infinitesimally-smooth}{无穷小光滑}（同上）。设 $S$ 是 artin \Term{local-ring}{局部环}的谱，$S_0$ 是由\Term{nilpotent-ideal}{幂零理想}定义的 $S$ 的\Term{subscheme}{子概形}，$T_0\in F_H(S_0)$；证明 $T_0$ \Term{liftable}{可提升}为 $F_H(S)$ 的元素 $T$。由假设（8.12 iv)），$T_0$ \Term{liftable}{可提升}为 $T'\in F(S)$。另一方面，由于 $H$ \Term{smooth}{光滑}，$T_0$ \Term{liftable}{可提升}为 $H_S$ 的\Term{subtorus}{子环面} $T''$（第 IX 讲 3.6 bis）；它与 $T'$ 由某个 $g\in G(S)$ \Term{conjugacy}{共轭}（同上），故 $T''\in F(S)$（8.12 v)）。由于 $\operatorname{Centr}_{G_S}(T'')$ 在 $S$ 上\Term{smooth}{光滑}，且在 $S_0$ 上与 $H_{S_0}$ 一致，$T''$ 位于 $H_S$ 的\Term{center}{中心}（第 IX 讲 5.6 a)），其在 $G$ 中的\Term{centralizer}{中心化子}等于 $H_S$；总之，$T''\in F_H(S)$，取 $T=T''$ 即可。

\begin{SGAresult}{推论}{8.13}
设 $S$ 是\Term{prescheme}{预概形}，$G$ 是在 $S$ 上\Term{smooth}{光滑}且\Term{finite-presentation}{有限呈示}、\Term{fiber}{纤维}\Term{connected}{连通}的 $S$-\Term{group-prescheme}{群预概形}。设 $\mathcal{TD}$ 是 $\mathcal T$ 的子\Term{functor}{函子}，其在 $S$-\Term{prescheme}{预概形} $S'$ 中取值的点集，是 $G_{S'}$ 中满足下述条件的\Term{subtorus}{子环面} $T$ 所成的\Term{set}{集合}：对 $S'$ 的每个点 $s'$，$T_{s'}$ 包含于 $G_{s'}$ 的\Term{derived-subgroup}{导出子群}（第 VI$_B$ 讲 7）中。
\SourcePage{319}
则 $\mathcal{TD}$ 可由在 $S$ 上\Term{smooth}{光滑}且\Term{separated}{分离}的 $S$-\Term{prescheme}{预概形}表示。
\end{SGAresult}
\begin{SGAresult}{推论}{8.13 bis}
设 $S,G$ 如上，$T$ 是 $S$-\Term{torus}{环面}，以
\[\operatorname{Hom}^{\mathrm{der}}_{S\text{-}\mathrm{gr}}(T,G)\]
表示 $\operatorname{Hom}_{S\text{-}\mathrm{gr}}(T,G)$ 的如下子\Term{functor}{函子}：其在 $S'$ 中取值的点集为满足以下条件的 $S'$-\Term{morphism}{态射} $u:T_{S'}\to G_{S'}$ 所成的\Term{set}{集合}：对 $S'$ 的每个点 $s'$，$u_{s'}$ 经 $G_{s'}$ 的\Term{derived-subgroup}{导出群}\Term{resolution}{分解}。则此\Term{functor}{函子}可由在 $S$ 上\Term{smooth}{光滑}且\Term{separated}{分离}的 $S$-\Term{prescheme}{预概形}表示。
\end{SGAresult}
推论 8.13 与 8.13 bis 的证明类似；例如证明 8.13 bis。用通常的方法归结到 $S$ 是 noether 的情形。核实 8.12 bis 的五条条件。

条件 i)、iv) 立即来自\Term{functor}{函子} $\operatorname{Hom}_{S\text{-}\mathrm{gr}}(T,G)$ 的相应性质。由于\Term{algebraic-group}{代数群}的\Term{derived-subgroup}{导出群}是\Term{normal-subgroup}{正规子群}（第 VI$_B$ 讲 \S7），条件 v) 成立。为证明 ii)，通过标准\Term{reduced}{约化}（EGA IV 8），只须证明：若 $S$ 是\Term{generic-point}{泛点}为 $\eta$ 的整 noether \Term{scheme}{概形}，$u:T\to G$ 是一个在泛\Term{fiber}{纤维}上经 $G_\eta$ 的\Term{derived-subgroup}{导出群}\Term{resolution}{分解}的 $S$-\Term{group-morphism}{群态射}，则存在 $\eta$ 的一个\Term{neighborhood}{邻域} $U$，使对 $U$ 的每个点 $s$，$u_s$ 经 $G_s$ 的\Term{derived-subgroup}{导出群}\Term{resolution}{分解}。这立即来自第 VI$_B$ 讲 10.12。

为证明 iii)，沿用 4.3 的记号。为证明
\[(u_m)_{m\in\mathbb N}\in\varprojlim_m\operatorname{Hom}^{\mathrm{der}}_{S_m\text{-}\mathrm{gr}}(T_m,G_m)\]
在 4.3 的意义下“容许”，且来自 $\operatorname{Hom}^{\mathrm{der}}_{S\text{-}\mathrm{gr}}(T,G)$ 的一个元素，立即归结到 $S$ 是\Term{complete}{完备}\Term{discrete-valuation-ring}{离散赋值环}的谱的情形（参见 8.1）。以 $t,s$ 分别表示 $S$ 的\Term{generic-point}{泛点}和\Term{closed-point}{闭点}，$D$ 表示 $G_t$ 的\Term{derived-subgroup}{导出群} $D_t$ 在 $G$ 中的\Term{schematic-closure}{概形闭包}。

a) $D_s$ 包含 $G_s$ 的\Term{derived-subgroup}{导出群}。事实上，由于 $D_t$ 在 $G_t$ 中\Term{normal}{正规}，且 $G$ 在 $S$ 上\Term{flat}{平坦}，$D$ 在 $G$ 中\Term{normal}{正规}。此外，\Term{morphism}{态射}
\[G\times_S G\longrightarrow G,\qquad(x,y)\longmapsto xyx^{-1}y^{-1}\]
在泛\Term{fiber}{纤维}上经 $D_t$ \Term{resolution}{分解}，故它经 $D$ \Term{resolution}{分解}。因此\Term{algebraic-group}{代数群} $G_s/D_s$ 交换，这就证明 a)。

b) 若 $u_m\in\operatorname{Hom}^{\mathrm{der}}_{S_m\text{-}\mathrm{gr}}(T_m,G_m)$，则 $u_m$ 经 $D_m$ \Term{resolution}{分解}。事实上，由假设 $u_s$ 经 $G_s$ 的\Term{derived-subgroup}{导出群}\Term{resolution}{分解}，故由 a)，更经 $D_s$ \Term{resolution}{分解}。由于 $D_m$ 在 $S_m$ 上\Term{flat}{平坦}，且在 $G_m$ 中\Term{normal}{正规}，\Term{quotient-group}{商群} $H_m=G_m/D_m$ \Term{representable}{可表}（第 VI$_A$ 讲 \S4）。由于 $T_m$ 在 $H_m$ 中的像是一个闭\Term{fiber}{纤维}为零的\Term{torus}{环面}（第 IX 讲 6.8），$T_m$ 的像为零，即 $u_m$ 经 $D_m$ \Term{resolution}{分解}。

c) 族 $(u_m)_{m\in\mathbb N}$“容许”，且\Term{liftable}{可提升}为属于 $\operatorname{Hom}^{\mathrm{der}}_{S\text{-}\mathrm{gr}}(T,G)$ 的\Term{morphism}{态射} $T\to G$。结合以上结论及 4.1 bis，只须证明 $D_t$ 是\Term{affine}{仿射}\Term{algebraic-group}{代数群}，这来自以下\Term{lemma}{引理}。
\begin{SGAresult}{\Term{lemma}{引理}}{8.14}
设 $G$ 是定义在域 $k$ 上的\Term{smooth}{光滑}\Term{connected}{连通}\Term{algebraic-group}{代数群}。则 $G$ 的\Term{derived-subgroup}{导出群} $D$ \Term{affine}{仿射}。
\end{SGAresult}
\SourcePage{320}
由于 $D$ 的形成与\Term{base-field-extension}{基域扩张}交换（第 VI$_B$ 讲 7），可以假定 $k$ \Term{algebraically-closed}{代数闭}。此时 $G$ 是由\Term{linear-group}{线性群} $L$ 对\Term{abelian-variety}{阿贝尔簇} $A$ 所作的\Term{extension}{扩张}。由于 $A$ 交换，$D$ 必包含于 $L$，故\Term{affine}{仿射}。

\textbf{\Term{maximal-torus}{极大环面}。}
\begin{SGAresult}{\Term{theorem}{定理}}{8.15}
设 $S$ 是\Term{locally-noetherian}{局部 noether}\Term{prescheme}{预概形}，$G$ 是在 $S$ 上\Term{smooth}{光滑}且\Term{finite-type}{有限型}的 $S$-\Term{group-prescheme}{群预概形}。则下列条件\Term{equivalence}{等价}：
\begin{enumerate}
\item[i)] $S$-\Term{functor}{函子} $\mathcal{TM}$ \Term{representable}{可表}，其中它在 $S$-\Term{prescheme}{预概形} $S'$ 中取值的点集等于 $G_{S'}$ 的\Term{maximal-torus}{极大环面}所成的\Term{set}{集合}（第 XII 讲 1.3）。
\item[ii)] 上述\Term{functor}{函子} $\mathcal{TM}$ 可由在 $S$ 上\Term{smooth}{光滑}且\Term{quasi-projective}{拟射影}、\Term{fiber}{纤维}\Term{affine}{仿射}的 $S$-\Term{prescheme}{预概形}表示。
\item[iii)] 群 $G$ 对\Term{etale-topology}{平展拓扑}\Term{local-scheme}{局部}有\Term{maximal-torus}{极大环面}。
\item[iv)] 群 $G$ 对\Term{faithfully-flat-topology}{忠实平坦拓扑}\Term{local-scheme}{局部}有\Term{maximal-torus}{极大环面}。
\item[v)] 群 $G$ 具有性质 AT（8.4），且其\Term{fiber}{纤维}的\Term{reductive-rank}{约化秩}是 $S$ 上的\Term{locally-constant}{局部常值}函数。
\end{enumerate}
\end{SGAresult}
\begin{proof}
$\mathrm{ii)}\Rightarrow\mathrm{i)}$ 显然。

$\mathrm{i)}\Rightarrow\mathrm{iii)}$。由于 $G$ 在 $S$ 上\Term{finite-presentation}{有限呈示}，由 6.3，$\mathcal{TM}$ 与环的\Term{inductive-limit}{归纳极限}交换，所以若它\Term{representable}{可表}，便是\Term{locally-finite-presentation}{局部有限呈示}的（EGA IV 8.14）。另一方面，$\mathcal{TM}$ \Term{formally-smooth}{形式光滑}（第 XI 讲 2.1 bis）。因此若它\Term{representable}{可表}，便由在 $S$ 上\Term{smooth}{光滑}的\Term{prescheme}{预概形}表示，再由 Hensel \Term{lemma}{引理}（第 XI 讲 1.10）得到 iii)。

$\mathrm{iii)}\Rightarrow\mathrm{iv)}$ 显然。

$\mathrm{iv)}\Rightarrow\mathrm{v)}$。设 $s\in S$。由假设，存在一个在 $S$ 上\Term{flat}{平坦}、像包含 $s$ 的 $S$-\Term{prescheme}{预概形} $S'$，使 $G_{S'}$ 有\Term{maximal-torus}{极大环面} $T'$。设 $s'\in S'$ 位于 $s$ 之上。于是 $G_{S'}$ 的\Term{fiber}{纤维}的\Term{reductive-rank}{约化秩}在 $\operatorname{Spec}\mathcal O_{S',s'}$ 上常值，因而 $G$ 的\Term{fiber}{纤维}的\Term{reductive-rank}{约化秩}在它的像 $\operatorname{Spec}\mathcal O_{S,s}$ 上常值（EGA IV 2.3.4 ii)）。而 $G$ 的\Term{fiber}{纤维}的\Term{reductive-rank}{约化秩}是 $S$ 上的\Term{local-scheme}{局部}\Term{constructible}{可构造}函数（6.3 bis），故此秩在 $s$ 的某个\Term{neighborhood}{邻域}上常值（EGA IV 1.10.1）。

为证明 $G$ 具有性质 AT，考虑一个 $S$-\Term{scheme}{概形} $S_1$，它是\Term{discrete-valuation-ring}{离散赋值环}的谱，且其\Term{closed-point}{闭点} $s_1$ 映到前述点 $s$。\Term{prescheme}{预概形} $S'_1=S_1\times_S S'$ 在 $S_1$ 上\Term{faithfully-flat}{忠实平坦}，而 $G_{S'_1}$ 有\Term{maximal-torus}{极大环面}。设 $A$（相应地 $A'$）为 $S_1$（相应地 $S'_1$）的环。将 $A'$ 视为其\Term{finite-type}{有限型}子 $A$-\Term{algebra}{代数}的\Term{inductive-limit}{归纳极限}，由 6.3，存在 $S_1$-\Term{scheme}{概形} $S''_1$，使 $G_{S''_1}$ 有\Term{maximal-torus}{极大环面}，且\Term{structural-morphism}{结构态射} $S''_1\to S_1$ \Term{finite-type}{有限型}且\Term{surjectivity}{满射}。再用 EGA II 7.1.9，可假定 $S'_1$ 是\Term{discrete-valuation-ring}{离散赋值环}的谱。此时显然 $G_{S'_1}$ 具有性质 AT，所以 $G_{S_1}$ 也具有此性质。由于对每个为\Term{discrete-valuation-ring}{离散赋值环}的谱的 $S$-\Term{prescheme}{预概形} $S_1$ 都如此，$G$ 具有性质 AT。

$\mathrm{v)}\Rightarrow\mathrm{i)}$。由 8.9，$G$ 的\Term{subtorus}{子环面}\Term{functor}{函子} $\mathcal T$ \Term{representable}{可表}；显然 $\mathcal{TM}$ 可由 $\mathcal T$ 中表示秩为 $r$ 的\Term{torus}{环面}所成的子\Term{functor}{函子}的既开又闭\Term{subprescheme}{子预概形}表示，
\SourcePage{321}
其中 $r$ 表示 $G$ 的\Term{reductive-rank}{约化秩}（可假定为常值）。

$\mathrm{iii)}\Rightarrow\mathrm{ii)}$。若条件 iii) 成立，便可使用第 XII 讲 7.1 的结果。于是\Term{functor}{函子} $\mathcal{TM}$ \Term{canonical}{典范}地\Term{isomorphism}{同构}于 $G$ 的\Term{cartan-subgroup}{Cartan 子群}\Term{functor}{函子}，只须应用 7.3 i)。
\end{proof}
\begin{SGAresult}{注记}{8.16}
可以证明，$G$ 的\Term{maximal-torus}{极大环面}\Term{prescheme}{预概形} $\mathcal{TM}$ 在 $S$ 上\Term{affine}{仿射}\footnote{参见 M. Raynaud，\emph{\Term{group-scheme}{群概形}与\Term{homogeneous-space}{齐性空间}上的\Term{ample-sheaf}{丰沛层}}（Faisceaux amples sur les schémas en groupes et les espaces homogènes），学位论文，待刊。（编者注：见 Lecture Notes in Mathematics 119（1970），Springer。）特别参见 IX 2.9。}，这推广第 XII 讲 5.4。
\end{SGAresult}
\begin{SGAresult}{推论}{8.17}
设 $S$ 是\Term{prescheme}{预概形}，$G$ 是在 $S$ 上\Term{smooth}{光滑}且\Term{finite-presentation}{有限呈示}的 $S$-\Term{group-prescheme}{群预概形}。假定 $G$ 的\Term{fiber}{纤维}的\Term{abelian-rank}{阿贝尔秩}与\Term{reductive-rank}{约化秩}是 $S$ 上的\Term{locally-constant}{局部常值}函数，则 $G$ 满足 8.15 中的（\Term{equivalence}{等价}）性质 i) 至 iv)。
\end{SGAresult}
可以假定 $G$ 的\Term{fiber}{纤维}的\Term{abelian-rank}{阿贝尔秩}与\Term{reductive-rank}{约化秩}均常值。仿照 8.11，结合 6.3 bis，归结到 $S$ 是 noether 的情形。此时 $G$ 具有性质 AT（8.7），再由 8.15 v) 得到结论。
\begin{SGAresult}{推论}{8.18}
设 $S$ 是\Term{prescheme}{预概形}，$G$ 是在 $S$ 上\Term{smooth}{光滑}的 $S$-\Term{group-prescheme}{群预概形}。则下列条件\Term{equivalence}{等价}：
\begin{enumerate}
\item[a)] $G$ 的\Term{fiber}{纤维}的\Term{unipotent-rank}{幺幂秩} $\rho_u$ 与\Term{abelian-rank}{阿贝尔秩} $\rho_{\mathrm{ab}}$（6.1 ter）是 $S$ 上的\Term{locally-constant}{局部常值}函数。
\item[b)] \Term{unipotent-rank}{幺幂秩} $\rho_u$ \Term{locally-constant}{局部常值}，且 $G$ 对 fpqc \Term{topology}{拓扑}\Term{local-scheme}{局部}有\Term{maximal-torus}{极大环面}。
\item[c)] $G$ 的\Term{fiber}{纤维}的\Term{reductive-rank}{约化秩} $\rho_r$（6.1 ter）与\Term{abelian-rank}{阿贝尔秩} $\rho_{\mathrm{ab}}$ 是 $S$ 上的\Term{locally-constant}{局部常值}函数。
\end{enumerate}
\end{SGAresult}
\begin{SGAresult}{注记}{8.19}
在 8.18 的假设下，用更精细的论证，并利用\Term{abelian-rank}{阿贝尔秩}的\Term{lower-semicontinuous}{下半连续}性（第 X 讲 8.7 中已宣布），可以证明：三个秩 $\rho_u,\rho_r,\rho_{\mathrm{ab}}$ 中若有两个\Term{locally-constant}{局部常值}，则第三个也\Term{locally-constant}{局部常值}。
\end{SGAresult}
\textbf{8.18 的证明。} 用 $G$ 的\Term{identity-component}{单位连通分支}代替 $G$ 后，可以假定 $G$ 在 $S$ 上\Term{finite-presentation}{有限呈示}（第 VI$_B$ 讲 5.3.3）。

$\mathrm{a)}\Rightarrow\mathrm{c)}$。设 $s\in S$。由于 $\rho_{\mathrm{ab}}$ \Term{locally-constant}{局部常值}，由 8.11，作一个\Term{covering}{覆盖} $S$ 的\Term{etale}{平展}\Term{group-extension}{扩张}后，可假定存在 $G$ 的\Term{subtorus}{子环面} $T$，其\Term{fiber}{纤维} $T_s$ 是 $G_s$ 的\Term{maximal-torus}{极大环面}。置 $C=\operatorname{Centr}_G(T)$，它是 $G$ 中在 $S$ 上\Term{smooth}{光滑}且\Term{fiber}{纤维}\Term{connected}{连通}的\Term{group-subprescheme}{群子预概形}。对 $S$ 的每一点 $t$，$C_t$ 显然包含 $G_t$ 的一个\Term{cartan-subgroup}{Cartan 子群} $C'_t$。缩小 $S$ 后，
\SourcePage{322}
可假定 $\rho_u,\rho_{\mathrm{ab}}$ 及 $C$ 在 $S$ 上的\Term{relative-dimension}{相对维数}均常值。于是有不等式
\[\begin{aligned}
\dim C_s=\dim C_t&\geq\dim C'_t\geq\rho_u(t)+\rho_{\mathrm{ab}}(t)+\dim T_t\\
&=\rho_u(s)+\rho_{\mathrm{ab}}(s)+\dim T_s=\rho_\nu(s)=\dim C_s.
\end{aligned}\]
可知 $C_t=C'_t$，所以 $T$ 是 $G$ 的\Term{maximal-torus}{极大环面}，从而更有 $\rho_r(t)=\rho_r(s)$。

$\mathrm{c)}\Rightarrow\mathrm{b)}$ 由 8.17 得到。

$\mathrm{b)}\Rightarrow\mathrm{a)}$。由于 $G$ 对 fpqc \Term{topology}{拓扑}\Term{local-scheme}{局部}有\Term{maximal-torus}{极大环面}，它也对 fpqc \Term{topology}{拓扑}\Term{local-scheme}{局部}有\Term{cartan-subgroup}{Cartan 子群}；从而（第 XII 讲 7.3），\Term{nilpotent-rank}{幂零秩} $\rho_\nu=\rho_u+\rho_r+\rho_{\mathrm{ab}}$ \Term{locally-constant}{局部常值}。由于 $\rho_r,\rho_{\mathrm{ab}}$ \Term{locally-constant}{局部常值}，$\rho_u$ \Term{locally-constant}{局部常值}。
